% Options for packages loaded elsewhere
\PassOptionsToPackage{unicode}{hyperref}
\PassOptionsToPackage{hyphens}{url}
\documentclass[
]{article}
\usepackage{xcolor}
\usepackage{amsmath,amssymb}
\setcounter{secnumdepth}{-\maxdimen} % remove section numbering
\usepackage{iftex}
\ifPDFTeX
  \usepackage[T1]{fontenc}
  \usepackage[utf8]{inputenc}
  \usepackage{textcomp} % provide euro and other symbols
\else % if luatex or xetex
  \usepackage{unicode-math} % this also loads fontspec
  \defaultfontfeatures{Scale=MatchLowercase}
  \defaultfontfeatures[\rmfamily]{Ligatures=TeX,Scale=1}
\fi
\usepackage{lmodern}
\ifPDFTeX\else
  % xetex/luatex font selection
\fi
% Use upquote if available, for straight quotes in verbatim environments
\IfFileExists{upquote.sty}{\usepackage{upquote}}{}
\IfFileExists{microtype.sty}{% use microtype if available
  \usepackage[]{microtype}
  \UseMicrotypeSet[protrusion]{basicmath} % disable protrusion for tt fonts
}{}
\makeatletter
\@ifundefined{KOMAClassName}{% if non-KOMA class
  \IfFileExists{parskip.sty}{%
    \usepackage{parskip}
  }{% else
    \setlength{\parindent}{0pt}
    \setlength{\parskip}{6pt plus 2pt minus 1pt}}
}{% if KOMA class
  \KOMAoptions{parskip=half}}
\makeatother
\usepackage{longtable,booktabs,array}
\usepackage{caption}
\captionsetup[table]{skip=6pt}
\newcounter{none} % for unnumbered tables
\usepackage{calc} % for calculating minipage widths
% Correct order of tables after \paragraph or \subparagraph
\usepackage{etoolbox}
\makeatletter
\patchcmd\longtable{\par}{\if@noskipsec\mbox{}\fi\par}{}{}
\makeatother
% Allow footnotes in longtable head/foot
\IfFileExists{footnotehyper.sty}{\usepackage{footnotehyper}}{\usepackage{footnote}}
\makesavenoteenv{longtable}
\usepackage{graphicx}
\makeatletter
\newsavebox\pandoc@box
\newcommand*\pandocbounded[1]{% scales image to fit in text height/width
  \sbox\pandoc@box{#1}%
  \Gscale@div\@tempa{\textheight}{\dimexpr\ht\pandoc@box+\dp\pandoc@box\relax}%
  \Gscale@div\@tempb{\linewidth}{\wd\pandoc@box}%
  \ifdim\@tempb\p@<\@tempa\p@\let\@tempa\@tempb\fi% select the smaller of both
  \ifdim\@tempa\p@<\p@\scalebox{\@tempa}{\usebox\pandoc@box}%
  \else\usebox{\pandoc@box}%
  \fi%
}
% Set default figure placement to htbp
\def\fps@figure{htbp}
\makeatother
\ifLuaTeX
  \usepackage{luacolor}
  \usepackage[soul]{lua-ul}
\else
  \usepackage{soul}
\fi
\setlength{\emergencystretch}{3em} % prevent overfull lines
\providecommand{\tightlist}{%
  \setlength{\itemsep}{0pt}\setlength{\parskip}{0pt}}
\ifPDFTeX
  \TeXXeTstate=1
  \newcommand{\RL}[1]{\beginR #1\endR}
  \newcommand{\LR}[1]{\beginL #1\endL}
  \newenvironment{RTL}{\beginR}{\endR}
  \newenvironment{LTR}{\beginL}{\endL}
\fi
\ifluatex
  \newcommand{\RL}[1]{\bgroup\textdir TRT#1\egroup}
  \newcommand{\LR}[1]{\bgroup\textdir TLT#1\egroup}
  \newenvironment{RTL}{\textdir TRT\pardir TRT\bodydir TRT}{}
  \newenvironment{LTR}{\textdir TLT\pardir TLT\bodydir TLT}{}
\fi
\usepackage{bookmark}
\IfFileExists{xurl.sty}{\usepackage{xurl}}{} % add URL line breaks if available
\urlstyle{same}
% fallback for those not using the hyperref driver hyperxmp:
\makeatletter
\@ifundefined{xmpquote}{\newcommand{\xmpquote}[1]{#1}}{}
\makeatother
\hypersetup{
  hidelinks,
  pdfcreator={LaTeX via pandoc}}

\author{}
\date{}

\begin{document}

\includegraphics[width=4.8394in,height=2.01563in]{media/image1.png}\LR{}

\LR{}

\LR{}

\LR{}

Thesis Proposal

Predicting Clinical Outcomes and Assessing the Effects of Diabetes Care
Guidelines Compliance in Hospitals Using Time-Dependent Machine Learning
and Deep Learning Methods

August 2025

\textbf{Submitted by:}

Shahar Oded

Department of Information Systems Engineering

Ben-Gurion University of the Negev

E-Mail:
\href{mailto:odedshah@post.bgu.ac.il}{\ul{odedshah@post.bgu.ac.il}}

\textbf{Supervised by:}

Prof. Yuval Shahar

Department of Information Systems Engineering

Ben-Gurion University of the Negev

E-Mail: \href{mailto:yshahar@bgu.ac.il}{\ul{yshahar@bgu.ac.il}}

Dr. Irit Hochberg

Institute of Endocrinology, Diabetes and Metabolism

Rambam Medical Center

E-Mail:
\href{mailto:i_hochberg@rambam.health.gov.il}{\ul{i\_hochberg@rambam.health.gov.il}}

\hyperref[abstract]{\textbf{Abstract 3}}

\hyperref[introduction]{\textbf{1 Introduction 4}}

\hyperref[related-work]{\textbf{2 Related Work 5}}

\begin{quote}
\hyperref[knowledge-based-temporal-abstraction]{2.1 Knowledge Based
Temporal Abstraction 6}

\hyperref[automated-support-to-guidelines-based-care]{2.2 Automated
Support to Guidelines Based Care 8}

\hyperref[assessing-the-quality-of-guideline-based-care]{2.3 Assessing
the Quality of Guideline-Based Care 10}

\hyperref[time-interval-mining]{2.4 Time Interval Mining 12}

\hyperref[sequential-models-in-medical-outcome-prediction]{2.5
Sequential Models in Medical Outcome Prediction 15}

\hyperref[explainable-ai-explaining-a-models-decision]{2.6 Explainable
AI: Explaining a Model's Decision 18}
\end{quote}

\hyperref[methods]{\textbf{3 Methods 20}}

\begin{quote}
\hyperref[computational-methods]{3.1 Computational Methods 21}

\hyperref[temporal-data-abstraction-and-qa-scores-calculation]{3.1.1
Temporal Data Abstraction and QA Scores Calculation 22}

\hyperref[predictive-model-development]{3.1.2 Predictive Model
Development 23}

\hyperref[frequency-based-temporal-patterns-vector-space-modeling-tpftirpvsm]{3.1.2.1
Frequency-Based Temporal-Patterns Vector-Space Modeling (TPF/TIRP+VSM)
23}

\hyperref[sequential-models]{3.1.2.2 Sequential Models 24}

\hyperref[comparison-and-research-goals]{3.1.2.3 Comparison and Research
Goals 29}

\hyperref[hyperparameter-optimization]{3.1.3 Hyperparameter Optimization
30}

\hyperref[evaluation-methods]{3.2 Evaluation Methods 31}

\hyperref[evaluation-of-the-knowledge-base]{3.2.1 Evaluation of the
Knowledge-Base 33}

\hyperref[evaluation-of-the-predictive-models]{3.2.2 Evaluation of the
Predictive Models 35}

\hyperref[feature-importance-analysis-explainability-and-protocol-refinement]{3.2.3
Feature Importance Analysis: Explainability and Protocol Refinement 38}

\hyperref[adherence-assessment-before-and-after-interventions]{3.2.4
Adherence Assessment Before and After Interventions 38}
\end{quote}

\hyperref[initial-results-and-schedule]{\textbf{4 Initial Results and
Schedule 39}}

\begin{quote}
\hyperref[initial-results]{4.1 Initial Results 39}

\hyperref[planned-evaluation-methods]{4.2 Planned Evaluation Methods 42}

\hyperref[proposed-schedule]{4.3 Proposed Schedule 43}
\end{quote}

\hyperref[references]{\textbf{References 45}}

\hyperref[appendices]{\textbf{Appendices 47}}

\begin{quote}
\hyperref[appendix-1-abstractions-quality-assessment]{Appendix 1:
Abstractions Quality Assessment 47}

\hyperref[appendix-2-automated-tak-generation-tool]{Appendix 2:
Automated TAK Generation Tool 54}
\end{quote}

\section{}\label{section}

\section{Abstract}\label{abstract}

Diabetes affects over\,7\,\% of Israelis yet consumes nearly a third of
hospital spending, mainly through complications triggered by in‑hospital
dysglycemia. Although new insulin protocols and quality indicators
exist, adherence is uneven and its impact unclear. We present an
explainable framework that predicts short‑term complications
(≤\,14\,days) from early EMR data. Raw, time-stamped patient data
recorded in the first K days, 2 ≤\,K\,\textless14, of each admission,
are converted into clinically meaningful interval-based abstractions,
and then into temporal patterns composed of frequent combinations of
such intervals. The patterns are exploited as predictive features.
Quality Assessment (QA) scores, on a {[}0, 1{]} scale, which assess the
compliance of the attending medical staff to established clinical
guidelines, are computed for each patient, relevant period, and
guideline component. These scores are used as features as well. QA
scores are also used to assess the effect of several interventions
intended to increase compliance to the guidelines. Patient demographics
are used as additional features. Two approaches are compared:\,(i) a
temporal pattern vector representation baseline, where Temporal
Probabilistic Profiles (TPF) turn symbolic patterns into fixed‑length
pattern distribution vectors; and\,(ii) a decoder‑only Transformer that
reads the same tokens with inter‑event gaps encoded by Time2Vec to learn
temporal rhythms. SHapley values reveal the most influential patterns,
highlighting critical protocol steps. We will also track adherence
before and after two national interventions - a diabetes e‑learning
module and a persistent‑hyperglycemia indicator - to situate the
predictions. The study aims to yield scalable risk‑stratification tools
and actionable guidance for improving inpatient diabetes care.

\section{1 Introduction}\label{introduction}

Diabetes mellitus is a global epidemic that claimed 3.4 \,million lives
in\,2024\,{[}25{]}. In Israel, people with diabetes already make up more
than\,7\,\% of the population and account for about one‑third of total
health‑care expenditure\,{[}26{]}. They occupy a disproportionate share
of inpatient days and drive costly complications such as infection,
acute kidney injury, cardiovascular events and diabetic
ketoacidosis\,{[}27{]}. Hyper‑ and hypoglycemia during admission extend
length of stay and raise mortality, yet adherence to evidence‑based
insulin protocols is inconsistent. Recent interventions - updated
guidelines, a diabetes e‑learning module and new quality indicators for
persistent hyperglycemia - aim to improve compliance, but their
real‑world impact remains unknown. Advanced analytics that quantify
guideline adherence at scale and predict short‑term complications are
therefore urgently needed.

Our study aims to predict the short term medical outcomes of diabetes
treatments (eg. Myocardial Infarction, Hypoglycemia, Ketoacidosis) by
leveraging Temporal Data Mining methods to process
patients\textquotesingle{} \emph{EMR} data, demographic information, and
treatment \emph{Quality Assessment (QA)} scores as features for short
term prediction models. If an Internal Ward admission is N days long,
we'll use information from the first K days to predict outcomes in the
subsequent N-K days (for 48 hours \textless= K \textless{} N \textless=
14 days admissions). We plan to evaluate multiple approaches, including:

\begin{enumerate}
\def\labelenumi{\arabic{enumi}.}
\item
  Using temporal intervals, gradients and patterns alongside demographic
  data as features;
\item
  Representing these patterns the distribution of frequent temporal
  patterns, i.e., the \emph{Temporal Probabilistic proFile (TPF)}
  method;
\item
  Utilizing \emph{Decoder-Only} \emph{Transformer} architecture capable
  of processing interval-based temporal abstractions and patterns (such
  as produced by the \emph{Temporal Mediator}) to generate a predicted
  events stream as basis for classification.
\end{enumerate}

A novel aspect of this research is the integration of medical knowledge
from two sources: general declarative knowledge (e.g., definitions such
as ``high blood pressure'') and procedural quality-control knowledge
(e.g., adherence to Insulin medical guidelines). Furthermore, by
employing explainable AI techniques like SHapley Additive ExPlanations
(SHAP), we aim to identify the most influential features in predicting
treatment outcomes, particularly the role of treatment adherence and
guideline compliance. This analysis will offer actionable insights into
clinical practices, highlighting which guidelines are critical,
redundant, or challenging to implement.

In addition, since the guideline-compliance scores are key features in
our predictive model, we will investigate how they may have been
influenced by recent interventions aimed at improving clinical care.
Specifically, we will assess the impact of two initiatives: (1) the
implementation of learnware software to train clinical staff on diabetes
management and (2) the introduction of a new quality measure for
managing persistent hyperglycemia lasting two or more days.
Understanding the effect of these interventions will help contextualize
the model's predictions and improve its relevance for evaluating care
quality and outcomes.

\section{2 Related Work}\label{related-work}

This chapter presents an overview of the key concepts and research areas
that serve as the foundation for this study. It explores existing
methodologies and systems in knowledge-based temporal abstraction,
automated support for guideline-based care, and the assessment of
guideline adherence - providing essential tools for integrating clinical
knowledge into predictive models.

We also examine advances in temporal data mining, focusing on sequential
pattern discovery from electronic medical records (EMRs) and the
integration of temporal patterns into predictive tasks. Additionally, we
introduce Transformer-based architectures, which have gained significant
traction in medical research, as an alternative to traditional data
mining methods. Finally, we highlight the growing importance of
explainable AI in clinical prediction, emphasizing methods that enhance
transparency in model decisions - critical for improving guideline
adherence and identifying key factors influencing treatment outcomes.

Each of these areas is discussed in detail in the following sections,
laying the groundwork for our study, which aims to predict short-term
medical outcomes through advanced temporal modeling and explainable
Machine Learning \& Deep Learning techniques.

\subsection{2.1 Knowledge Based Temporal
Abstraction}\label{knowledge-based-temporal-abstraction}

\textbf{The Knowledge-Based-Temporal-Abstraction (\emph{KBTA)}
principles} {[}24, 4{]} - A methodology that structures the task of
transforming time-stamped clinical data into abstract, interpretable
context dependent concepts (e.g., "decreasing temperature for 3 days").
It abstracts raw time-stamped data into interval-based meaningful
representations like states, events and gradients. It does so by
combining domain-specific knowledge of relationships and constraints,
data, and context from the patient's \emph{Electronic Health Record}.

\begin{itemize}
\item
  State abstractions: Representing raw numeric values (e.g., heart rate)
  as symbolic states (e.g., low, normal, high) over specific time
  intervals. Requires pre definition of the discretization thresholds
  and relevance duration (e.g., glucose measurement is relevant for up
  to 6 hours).
\item
  Gradient abstractions: Describing trends or changes in a parameter
  over time (e.g., increasing, stable, decreasing). Requires pre
  definition of what a meaningful change in a parameter is and over what
  duration (e.g., an increase of 5 {[}unit{]} over 12 hours).
\item
  Context-sensitive abstractions: Allowing for the generation of
  abstractions based on specific clinical contexts or conditions. For
  example, in the context of diabetes management, blood glucose levels
  can be abstracted differently to low, normal and high ranges depending
  on background diseases.
\end{itemize}

The representation of \emph{KBTA} is done using \emph{\textbf{TAK}}
language which defines the abstraction guidelines using different types
of logic files:

\begin{itemize}
\item
  \textbf{Raw concepts}: Basic clinical parameters, such as systolic
  blood pressure.
\item
  \textbf{Events}: An event in time, with no specific duration. For
  example, admission of a medication or a Ketoacidosis event.
\item
  \textbf{Context}: A temporal information that may affect the
  conditions applied when calculating the intervals, like Alcohol Abuse
  or prior surgery.
\item
  \textbf{State patterns}: High-level temporal patterns derived from raw
  concepts, such as repeated instances of elevated glucose levels over a
  specified period.
\item
  \textbf{Trend}: Represent trends (gradient) or changes in clinical
  parameters over time, such as increasing, stable, or decreasing blood
  glucose levels.
\item
  \textbf{Pattern}: Represent a sequence of temporal events/ intervals
  that occur at a specific order, for example Meal \ul{before} Insulin
  Admission \ul{before} Glucose level decreases.
\item
  \textbf{Quality Assessment (QA)}: Refers to compliance with clinical
  objectives or guidelines, coded as \emph{patterns}, such as
  administering insulin within a prescribed time frame in relation to
  meals. This concept expands on the original architecture, as detailed
  in Paper 2 by Shalom, Erez, et al., and will be further elaborated in
  \textbf{Section 2.3}.
\end{itemize}

\textbf{The Temporal Mediator} {[}4{]} serves as a system designed to
enable temporal reasoning and abstraction in clinical decision-making.
By automatically executing TAK files, it processes time-stamped clinical
data, such as blood test results, medical events (e.g., insulin
injections), and abstraction goals like pattern recognition or
therapeutic decision-making. The Mediator generates interval-based,
context-aware abstractions, such as "five days of increasing glucose
levels," to aid in diagnosis or treatment planning.

\includegraphics[width=6.58854in,height=2.8928in]{media/image24.png}

\textbf{Figure 1:} An illustration of the temporal abstraction process:
Multivariate clinical time series data is transformed into
interval-based abstractions representing states and gradients.

Unlike baseline temporal discretization methods-such as
\textbf{equal-length} (dividing the timeline into fixed-size intervals)
or \textbf{equal-frequency} (splitting the timeline so each bin contains
a similar number of observations)-KBTA does not impose arbitrary,
data-agnostic boundaries. Instead, it uses clinically meaningful
thresholds, relevance durations, and contextual rules to define
intervals that reflect the \emph{semantics} of the patient's state. This
knowledge-driven approach preserves the temporal coherence of clinical
events, captures durations as first-class features, and produces
abstractions that are both interpretable to clinicians and aligned with
real-world treatment logic, rather than being artifacts of uniform
binning strategies.

\subsection{2.2 Automated Support to Guidelines Based
Care}\label{automated-support-to-guidelines-based-care}

\textbf{Clinical Guidelines (GLs)} {[}2{]} are evidence-based protocols
developed to standardize medical care, improve treatment outcomes, and
reduce variability in clinical decisions. They aim to enhance care
quality, reduce errors, and improve patient outcomes by offering
structured decision pathways. They can make treatments more uniform,
improve decision quality, enhance patient survival, reduce morbidity,
and potentially lower costs. However, their implementation in practice
is complicated, as \emph{GLs} are text-based and not easily accessible
to therapists who need to adapt them to their patients and apply them at
the point of care, leading to low adherence and adaptation among
physicians.

Consequently, there is a need for \emph{Clinical Decision Support
Systems (CDSSs, healthcare specific DSSs)} that automatically apply and
assess \emph{GLs}, improving care quality.

This causes for an important task which is acquiring these \emph{GLs} in
a machine-interpretable manner that can be enforced algorithmically,
creating a \emph{medical knowledge base (KB)}.

\emph{\textbf{Asbru} {[}23, 2, 3{]}}- A computer-interpretable language
for \emph{GL} specification. Guideline representation language that
allows guideline designers to represent sequential, concurrent,
conditional and repetitive procedural actions as a hierarchy of
execution plans.

\emph{\textbf{DeGeL}} {[}2{]}- Digital Electronic Guidelines Library is
a web-based, distributed architecture, developed at BGU Medical
Informatics Research Center, that facilitates the gradual conversion of
clinical guidelines from text to a formal representation in
\emph{Asbru.} The conversion is handled through a graphical tool called
\emph{Gesher}, which helps in the construction process of the
\emph{GLs}, with the knowledge engineer helps in its formal
representation.

When transporting \emph{GLs} from textual instructions to
machine-readable, we need to formalize two types of medical knowledge
{[}2{]}:

\begin{itemize}
\item
  \emph{\textbf{Declarative Knowledge ("What Is"):}} Describes medical
  terms and conditions, such as definitions of diseases and diagnostic
  criteria (based on \emph{KBTA} principles). Utilized for understanding
  the underlying facts. This knowledge is represented in the \emph{TAK}
  language and is the abstraction guidelines for the \emph{Mediator}.
\item
  \emph{\textbf{Procedural Knowledge ("How To"):}} Specifies the
  sequence of actions, such as treatment steps and administration
  protocols. Represented in \emph{Asbru.}
\end{itemize}

\emph{\textbf{Picard}} is a continuous-care guideline application
framework and clinical decision support system (CDSS) designed to help
clinicians apply clinical guidelines over extended periods. It offers
real-time, data-driven recommendations based on
patients\textquotesingle{} longitudinal clinical data, with the goal of
improving guideline adherence, enhancing decision completeness, and
reducing errors and redundant actions.

The engine operates using several key components:

\begin{itemize}
\item
  Declarative Knowledge: Based on predefined \emph{TAK} files, the
  engine uses the \emph{Temporal Mediator} for temporal reasoning and to
  answer specific queries about patient data.
\item
  Procedural Knowledge: Guides the expected course of action based on
  clinical guidelines (\emph{GLs}).
\item
  \emph{DeGeL} Digital Library: Supplies procedural and declarative
  knowledge for clinical guidelines in a machine-readable format.
\item
  Service-Oriented Architecture (SOA): Enables modular, distributed
  deployment for clinicians, patients, and knowledge engineers
\end{itemize}

The engine processes time-stamped clinical data (e.g., lab results,
vital signs) and patient-specific context from electronic medical
records (EMRs) to generate actionable recommendations (e.g., "Administer
magnesium sulfate"), reports on guideline adherence, and alerts for
critical deviations or omissions {[}5,6{]}.

\emph{Picard} served as the core DSS engine in the \textbf{MobiGuide}
project, which supported chronic disease management through distributed
decision support (local mobile DSS and backend DSS) in Spain and Italy.
The engine helped improve clinicians\textquotesingle{} compliance with
guidelines dramatically, from 41\% in manual mode to 93\% with Picard.
The system provided real-time, patient-specific recommendations by
integrating data from hospital medical records, mobile biosensors, and
patient inputs. This approach enabled continuous monitoring and
personalized guidance {[}6,7{]}.

\emph{\textbf{e-Picard}} is an adaptation of \emph{Picard} for
non-continuous, episodic use where care is sought sporadically and not
in real-time. It conducts retrospective \emph{Quality Assessment} of
past actions using \emph{Temporal Fuzzy Logic} {[}\textbf{Section
2.3}{]}, and returns recommendations for future actions based on partial
or omitted implementations, all while incorporating patient-specific
contexts such as lifestyle or situational changes {[}1{]}. Such system
is usually evaluated on 2 criterias {[}6{]}:

\begin{itemize}
\item
  \textbf{correctness -} accuracy of recommendations.
\item
  \textbf{completeness} - coverage of guideline-recommended actions.
\end{itemize}

These are usually applied in 2 evaluation phases:

\begin{itemize}
\item
  \textbf{Technical Evaluation:} Focuses on ensuring the system
  generates the correct recommendations and avoids incorrect ones in
  predefined scenarios. It verifies the system\textquotesingle s
  functionality against a controlled, simulated longitudinal dataset
  tailored to predefined scenarios.
\item
  \textbf{Silent Evaluation:} Assesses the system\textquotesingle s
  performance on real-world data without presenting its recommendations
  to the treatment team, allowing for independent expert review of the
  system\textquotesingle s outputs.
\end{itemize}

\subsection{2.3 Assessing the Quality of Guideline-Based
Care}\label{assessing-the-quality-of-guideline-based-care}

\textbf{Fuzzy Logic} {[}2{]} - Serves as a nuanced approach to assessing
guideline adherence by assigning degrees of compliance rather than
binary outcomes. This method is particularly valuable in clinical
scenarios where actions may partially fulfill the ideal protocol. Using
\emph{Temporal Fuzzy Logic functions}, it evaluates adherence to
constraints such as timing, frequency, and value ranges of medical
actions. Scores range from 0 (non-compliance) to 1 (full compliance),
allowing for a more flexible and realistic assessment of care quality.

For example, a test required within 72--120 hours after admission would
score 1 if performed within 72 hours, with scores decreasing
proportionally for delays up to 120 hours. \emph{Fuzzy Logic} is applied
within \emph{Quality-Assessment Temporal Patterns (QATPs)} to evaluate
partial compliance with temporal constraints, such as periodic
assessments or gaps between linked actions. This approach is integral to
the \emph{Quality Assessment System}, enabling a detailed and automated
evaluation of healthcare interventions.

To retrospectively assess the execution of procedural knowledge within
patient data, as in the \emph{e-Picard} engine, a third type of
knowledge was developed in addition to the \emph{Declarative and
Procedural Knowledge}. This knowledge transforms the procedural aspects
of \emph{GLs} into declarative temporal-abstraction patterns,
representing the expected execution trace in patient records for correct
therapy application {[}2{]}.

\begin{itemize}
\item
  \emph{\textbf{Quality-Assessment (QA) Knowledge:}} Focuses on
  evaluating the intended outcomes and processes of guideline
  application. It facilitates the computation of quality-assessment
  scores by identifying \emph{Quality Assessment Temporal Patterns
  (QATPs)} in a patient's longitudinal record. These patterns serve as
  an "execution trace," indicating whether the guideline was implemented
  correctly. In fact, these files are written to represent the
  \emph{procedural knowledge} in a \emph{declarative} manner.
  Represented in the \emph{TAK} language with \emph{Temporal Fuzzy
  Logic} extensions, this knowledge supports partial evaluations of
  guideline adherence over time.
\end{itemize}

When acquiring the QA knowledge, a few steps are needed:

\begin{enumerate}
\def\labelenumi{\arabic{enumi}.}
\item
  Actions from the procedural guideline are categorized based on their
  attributes in Asbru's plan types (e.g., single-action, cyclical,
  sequential) and mapped to specific constraint types, which define how
  compliance is assessed. The key constraint types include:
\end{enumerate}

\begin{itemize}
\item
  Cyclical (Periodic): Specifies actions that need to be performed
  periodically, such as routine check-ups or recurring tests within
  defined intervals.
\item
  Temporal Gap: Represents the required time interval between
  consecutive actions or plans, ensuring appropriate spacing (e.g.,
  waiting a set number of hours between medication doses).
\item
  Temporal Relation: Captures relationships between actions, such as one
  action needing to precede or overlap with another (e.g., administering
  insulin before meals). The possible values are relations: precedes,
  meets, overlaps, starts, during, finishes, equal, with the default
  being precedes.
\item
  Value: Assesses compliance based on specific thresholds or ranges for
  measured values, such as ensuring a patient\textquotesingle s glucose
  level remains within an acceptable range during treatment.
\end{itemize}

\begin{quote}
These constraint types ensure that procedural actions are systematically
mapped and compliance is evaluated with precision, reflecting the
intended execution of the guideline. Cyclical (Periodic), Temporal Gap,
and Value constraints utilize \emph{Temporal Fuzzy Logic} to enable
partial compliance assessments. Fuzziness accounts for variability in
frequency for cyclical actions, time gaps between consecutive actions,
and measured values, allowing for flexible and realistic evaluations of
adherence to guidelines.
\end{quote}

\begin{enumerate}
\def\labelenumi{\arabic{enumi}.}
\setcounter{enumi}{1}
\item
  Defining \emph{QATPs} (Quality-Assessment Temporal Patterns):
  \emph{Temporal Fuzzy Logic} is applied to define \emph{QATPs}, which
  specify the constraints on actions and the calculation of partial
  compliance scores. This ensures nuanced evaluation by accounting for
  degrees of adherence. This is done using a trapezoidal model, in which
  the non-parallel edges represent the slopes on which a partial action
  will be scored, depending on its position on the slope.
\item
  Creating a Weighted QA Tree: The \emph{Procedural} \emph{GL} is
  translated into a weighted tree structure based on Asbru's hierarchy,
  where each node represents a plan with a defined weight. This tree
  serves as a compliance framework, linking procedural elements to their
  QA representation.
\item
  Compliance Assessment: Compliance is calculated recursively through
  the QA tree. Child nodes' compliance scores are aggregated, using
  weights, to derive the compliance score of their parent nodes. This
  hierarchical approach ensures that the overall adherence reflects both
  individual actions and their relationships.
\end{enumerate}

With the QATP extensions, the \emph{Mediator}, described in
\textbf{Section 2.1}, can retrospectively calculate adherence scores to
the guidelines, leveraging protocol data already transformed into
machine-readable files {[}2{]}. These scores will be used as additional
features in our study and its influence on the treatment's outcome will
be examined.

\subsection{2.4 Time Interval Mining}\label{time-interval-mining}

This chapter outlines the key methodologies for transforming
longitudinal electronic medical record (EMR) data into feature vector
representations that capture a patient's state up to a specific point in
time. We explore tools and techniques in temporal pattern mining that
form the foundation of our prediction models, enabling the conversion of
raw time-stamped data into structured and semantically rich features.

\textbf{KarmaLego} {[}9{]} is a framework that extends the
\emph{Knowledge-Based Temporal Abstraction (KBTA)} method and integrates
time-interval mining capabilities to identify Frequent \emph{Time
Interval Relation Patterns (TIRPs)} in multivariate time series. For
example, it can detect patterns like "a period of high-dose insulin
administration" preceding "two or more hypoglycemia events" using
Allen's seven basic interval relations. It identifies temporal
relationships such as "Before", "Meets'', or "Overlaps" when they appear
in more than p\% (minimum support) of the longitudinal patient records
(abstracted states by the \emph{Mediator}). The output are the frequent
patterns discovered in the data based on parameters.

KarmaLego begins where KBTA ends. KBTA generates symbolic intervals and
establishes the abstraction layer, while KarmaLego uses these intervals
as its input to discover and analyze complex temporal patterns across
multiple continuous features.

It was shown that both temporal abstractions of raw data and
combinations of these abstractions, which create more complex temporal
patterns, can serve as effective features for machine learning tasks
such as clustering, classification, prediction, visualization, and
interactive analysis of large temporal datasets.

The \emph{Semantic Adjacency Criterion (SAC)} {[}11{]} is a filtering
principle designed to enhance temporal pattern discovery by ensuring
that mined patterns are semantically coherent and domain-relevant. It
operates by applying semantic constraints during the mining process to
exclude patterns where intermediate intervals disrupt or contradict
meaningful relationships. By focusing on semantically valid patterns,
SAC reduces computational overhead significantly, simplifies datasets,
and improves the interpretability of results.

When integrated with KarmaLego, SAC can refine its output by ensuring
that the discovered \emph{Time Interval-Relation Patterns (TIRPs)} are
not only frequent but also align with domain-specific semantics.

An interesting usage example can be found in the study by Raymond,
Michal Weisman, and Yuval Shahar {[}12{]} which leveraged the
\emph{KBTA} and \emph{KarmaLego} tools to predict the next clinical
action for ICU patients. The study aimed to assist clinicians by
predicting the most likely next action based on patterns observed in
historical data of similar patients, providing valuable decision
support.

Using the \emph{MIMIC-IV} dataset - a comprehensive repository of
critical care admissions - the study focused on subsets of patients
diagnosed with Hypoglycemia, Hypokalemia, and Hypotension. It employed a
two-step machine learning approach: a Random Forest classifier was used
to determine whether a treatment was necessary, followed by a Gradient
Boosting Regressor to predict the appropriate dosage when a treatment
was deemed necessary.

This research highlights the potential for decision support systems to
enhance clinical practice by enabling medical teams to make more
informed and timely decisions regarding ICU patients. It also
demonstrates the utility of the data abstraction tools (\emph{KBTA} and
\emph{KarmaLego}) in transforming raw, time-stamped clinical data into
actionable insights. Moreover, the goal of predicting the next action
for a patient assumes the existence of an optimal ``best action'',
aligning with our methodology of using QA scores to predict treatment
outcomes.

\textbf{Dynamic Time Warping (DTW)} {[}18{]}: is a well-known method for
measuring similarity between time-series data of different lengths. It
works by finding an optimal alignment (or "warping") between two
sequences to minimize the distance between them. Originally developed
for speech recognition, it is now widely applied in various fields,
including medicine. DTW is effective in handling temporal misalignments
and identifying similar patterns even if they occur at different times.
However, classical DTW is limited to univariate sequences, assumes
complete datasets, and does not consider clinical abstractions.

\textbf{Interval-Based Dynamic Time Warping (iDTW)} {[}18{]}: extends
DTW to handle interval-based, multivariate, and abstracted data from
longitudinal medical records (e.g., "HIGH glucose for 2 days"). It
produces a distance measure quantifying the similarity between patient
records, enabling its use in machine learning. The method focuses on
differences in patients\textquotesingle{} longitudinal data and
introduces two key components:

\begin{itemize}
\item
  \emph{iRep (Interval Representation):} Transforms raw time-stamped
  clinical data into interval-based temporal abstractions using
  domain-specific knowledge (e.g., Knowledge-Based Temporal Abstraction
  - KBTA). Produces a 2D matrix of temporal abstractions that represent
  states (e.g., HIGH/LOW) and trends (e.g., INCREASING/DECREASING) for
  multiple variables at predefined granularities.
\item
  \emph{iMatch (Interval Matching):} Extends DTW by matching the
  multivariate interval-based representations. It calculates the
  distance between intervals using domain-specific metrics, considering
  both value and duration.
\end{itemize}

The paper further describes how this method was utilized in order to
create a distance matrix between patients, based on their longitudinal
data, which was then utilized by a KNN-like model to evaluate similarity
and perform classification.

\textbf{Temporal Probabilistic ProFiles (TPF)} {[}10{]} Is used for
classification and prediction tasks of multivariate time series. TPFs
are the distributions of the number (cardinality) or duration of the
frequent temporal patterns found in each patient's record (and by
averaging, in the whole patient's population). These profiles
incorporate temporal dependencies and uncertainties, allowing for
comparison between different profiles or used as meta-features for
classification tasks on the time series / outcome prediction (of events
such as medical conditions) by measuring their distance, either to the
aggregated vector of each class (\emph{the Category method}), or to the
individual profiles of each of the entities, using negative cross
entropy to find the closest K (\emph{the Top K method}). Unlike
TIRP-based classification, which relies on a pattern's horizontal
support as a classification feature, this method goes further by using
the pattern's relative horizontal support within a patient- compared to
other patients - or its duration, as indicators of its potential
importance in the treatment trajectory.

How it works:

\begin{itemize}
\item
  Input: Multivariate longitudinal data.
\end{itemize}

\begin{itemize}
\item
  Process:

  \begin{itemize}
  \item
    The data is abstracted to shorter interval based representation
    using the \emph{KBTA} principles (\emph{Mediator}).
  \item
    The data is mined for frequent patterns across groups of patients
    using the \emph{KarmaLego} algorithm.
  \item
    \emph{Temporal Probabilistic Profiles (TPFs)} are created based on
    the distribution of said frequent patterns, as a singular vector per
    patient up to time t.
  \end{itemize}
\item
  Output: A summary vector (profile) defining the important patterns
  (and their distribution) in a patient's EHR.
\end{itemize}

These profiles can be directly applied for classification using
vector-space methods like KNN or by measuring similarity to the
aggregated vector of the closest group. Additionally, they may serve as
features for more advanced classification models in our research.

\includegraphics[width=5.96354in,height=2.22917in]{media/image12.png}

\textbf{Figure 2:} Overview of the temporal-probabilistic profile-based
and prediction method.

In the same paper, the authors used this methodology to predict a sepsis
complication in ICU patients by integrating temporal data mining with
machine learning.

Using the \emph{MIMIC III} dataset, which includes 40,000 critical care
admissions from 2001 to 2012, this study demonstrated the effectiveness
of the \emph{TPF} abstraction method. The performance of a \emph{Random
Forest} model trained on \emph{TPF}-generated vectors was compared to
that of an \emph{RNN (GRU-D)} model trained on the original multivariate
raw data combined with patients' demographic information. The GRU-D
input is the event token and the time since the last observation of this
variable (the time gap/delta)

In various configurations, the \emph{TPF}-based approach paired with a
traditional machine learning model consistently outperformed the more
complex sequential model.

Despite dealing with a different dataset on a different prediction
problem, this proves the potential effectiveness of the \emph{TPF}
methodology, which might be of use to us in the representation of the
first K days before prediction.

\subsection{2.5 Sequential Models in Medical Outcome
Prediction}\label{sequential-models-in-medical-outcome-prediction}

Sequential models, particularly \emph{Transformers}, as introduced in
"Attention is All You Need" {[}13{]}, have emerged as a powerful
approach for tasks like outcome prediction by capturing dependencies
across temporal sequences. These Encoder-Decoder models consist of an
encoder that processes input sequences (e.g., time-series data
representing clinical states) and a decoder that generates output
sequences (e.g., predicted treatment actions or dosages). The attention
mechanism, a key innovation in this architecture, enables the model to
focus on the most relevant parts of the input sequence, addressing the
limitations of traditional recurrent models in handling long-term
dependencies. This mechanism is capable of capturing implicit patterns
in the temporal data, unlike \emph{KarmaLego} based methods that capture
explicit patterns.

Building on this foundation, I reviewed several studies that applied
sequential models - particularly machine learning and Transformer-based
approaches - on raw Electronic Health Records (\emph{EHRs}) for various
predictive tasks. I also reviewed experimental methods for embeddings
and positional encodings needed for the prediction task using models
such as these. This review provided valuable insights into preprocessing
strategies, feature engineering, and modeling approaches relevant to
structuring the experiments in this research.

The paper \emph{``\hl{Treatment Prediction in the ICU Using a
Partitioned, Sequential, Deep Time Series Analysis}'' {[}20{]}}
introduced a framework for treatment prediction in ICU settings using a
bidirectional LSTM model, leveraging multivariate time-series data
(e.g., vital signs, lab results) and static features (e.g.,
demographics) from the MIMIC-IV ICU database. The approach incorporated
partitioned prediction windows, splitting a 12-hour horizon into three
4-hour sub-windows, using data from the previous 24 hours for
predictions. Temporal features were imputed for hourly granularity,
while constant features were integrated into the LSTM, with internal
integration yielding better performance.

The prediction tasks focused on hypoglycemia (IV glucose), hypokalemia
(IV and oral potassium), and hypotension (IV fluids, Dopamine, and
Norepinephrine). The two-step model first predicted treatment necessity
and then estimated dosage. Sequential reuse of earlier sub-window data
enhanced training for later predictions.

This research demonstrated the effectiveness of temporal neural networks
in predicting both treatment necessity and dosage, mimicking ICU
decision-making processes. The integration of constant features and
partitioned predictions highlights its potential for decision support,
reducing clinician cognitive load while improving accuracy. Unlike our
study, this study does not mention the usage of abstraction methods or
QA scores, which might yield good results when incorporated with
bidirectional models such as the one mentioned here.

Moving beyond recurrent models, attention-based architectures such as
Transformers have further advanced this field. The paper
\emph{``TransformEHR: Transformer-Based Encoder-Decoder Model for
Disease Outcome Prediction''} {[}15{]} presents a generative
encoder-decoder architecture specifically designed for predicting future
disease outcomes using longitudinal EHR data. TransformEHR processes
multi-level embeddings of patient data, including ICD-10 diagnostic
codes, visit-level information, and demographic features. To incorporate
temporal information, the authors explicitly encode time gaps between
visits as sinusoidal embeddings of visit dates or day differences,
enabling the model to account for irregular visit intervals in clinical
practice. The embeddings for visit, demographic info and time were
summed up to a full visit embedding vector, fed into the transformer
encoder as a one-channel stream.

Unlike prior EHR-focused models (which use only encoder-based
architectures), TransformEHR's encoder-decoder transformer leverages
cross-attention to connect historical patient data to future disease
predictions, and self-attention in the decoder to sequentially predict
all outcomes within a future visit. This architecture was pre-trained to
recover all ICD codes of a future visit from past visits, a strategy
that outperformed conventional code-level masking used in BERT-style
models.

This framework demonstrated state-of-the-art performance in predicting
both common and rare conditions, including tasks like pancreatic cancer
onset prediction and intentional self-harm in PTSD patients. The study
also highlights challenges similar to ours, particularly in handling
multi-dimensional sequential EHR data and determining the optimal way to
encode time and integrate visit-level information into the transformer
architecture.

A third interesting paper I would like to mention is "Zero shot health
trajectory prediction using transformer" {[}14{]}. The paper introduces
\emph{ETHOS} (Enhanced Transformer for Health Outcome Simulation), a
decoder-only transformer architecture designed for zero-shot prediction
of Patient Health Timelines (PHTs). \emph{ETHOS} treats health events,
including admissions, lab tests, and medications, as tokens - with
static patient data (like demographics) embedded at the start of each
sequence - and feeds these flattened chronological sequences into the
transformer. Temporal gaps are encoded using explicit interval tokens
(e.g., 5-15 minutes, 1-6 months), rather than learned sinusoidal or
positional embeddings that feed on the exact 𝚫t. This approach enables
ETHOS to operate directly on irregular, noisy data, using quantile-based
tokens to mitigate bias from missing values.

\emph{ETHOS} demonstrates robust zero-shot performance across multiple
tasks, including ICU mortality prediction (AUC 0.927), hospital
mortality (AUC 0.921), length of stay estimation (MAE 2.26 days), and
even DRG classification (top-1 accuracy 84.8\%). Crucially, ETHOS
achieves these results without fine-tuning, highlighting the power of
unified foundation models for diverse clinical prediction tasks. The
study underscores the feasibility of deploying transformer-based models
trained once to perform a broad range of patient trajectory predictions,
addressing challenges in handling multidimensional, sequential, and
temporally irregular EHR data.

While both \emph{TransformEHR} and \emph{ETHOS} address temporal
dependencies in \emph{EHR} data, neither directly incorporates data
abstraction techniques (like \emph{Mediator}), pattern mining
(\emph{KarmaLego}), or \emph{QA scores} - methods we believe are crucial
for mitigating overfitting and highlighting clinically significant
intervals. This gap motivated us to explore temporal embeddings and
tokenization methods that can generalize across data granularities.

A solution for the gap in Δt representation was found in the paper
"Time2vec: Learning a vector representation of time" {[}19{]}. Time2Vec
is a general-purpose, model-agnostic representation of time that
addresses the challenge of incorporating temporal information in
sequence models. Unlike approaches that treat time as an ordinary
feature or rely on handcrafted time representations, Time2Vec learns a
compact vector embedding for each timestamp, capturing both periodic and
non-periodic patterns in the data. It consists of a linear term and a
series of periodic sine terms with learnable frequencies and phase
shifts. This enables it to encode complex temporal behaviors, such as
seasonality or cyclicity, in a data-driven manner. Crucially, Time2Vec
is invariant to time rescaling, ensuring robust performance across
different temporal granularities. By directly concatenating these
learned time embeddings with other input features, Time2Vec can be
seamlessly integrated into various models (like LSTMs or Transformers),
consistently improving performance across diverse datasets, as
demonstrated in the authors' experiments.

Building upon these insights, our sequential model leverages a
decoder-only transformer architecture enhanced with Time2Vec embeddings
and multiple levels of temporal abstraction to address the unique
challenges of modeling EMR data, and more specifically, address the
unique challenges of modelling abstracted EMR data (Mediator's output),
under the assumption that this more knowledgeable input will assist the
model in generating better predictions.

\subsection{2.6 Explainable AI: Explaining a Model's
Decision}\label{explainable-ai-explaining-a-models-decision}

The \textbf{SHAP} {[}16{]} framework (SHapley Additive ExPlanations)
provides a unified approach for interpreting machine learning
predictions by assigning importance values to individual features for a
specific prediction. Derived from cooperative game theory, \emph{SHAP}
values represent the marginal contribution of each feature to the
prediction, ensuring that these attributions adhere to desirable
properties like local accuracy, missingness, and consistency.

\begin{enumerate}
\def\labelenumi{\arabic{enumi}.}
\item
  Explanation Model:
\end{enumerate}

\begin{quote}
\emph{SHAP} represents a model's explanation as a linear function of
simplified inputs, where each feature's contribution is computed
relative to its impact on the output.
\end{quote}

\begin{enumerate}
\def\labelenumi{\arabic{enumi}.}
\setcounter{enumi}{1}
\item
  Adapting to Classical Machine Learning Models:
\end{enumerate}

\begin{quote}
For linear models, \emph{SHAP} values are derived directly from the
model coefficients, weighted by the deviation of feature values from the
dataset mean. For ensemble models (e.g., Random Forest, Gradient
Boosting), \emph{SHAP} uses game-theoretic approximations to efficiently
compute contributions.
\end{quote}

\begin{enumerate}
\def\labelenumi{\arabic{enumi}.}
\setcounter{enumi}{2}
\item
  Adapting to Seq2Seq Models:
\end{enumerate}

\begin{quote}
\emph{SHAP} can be extended beyond tabular settings to sequential
architectures, including RNNs, LSTMs, and transformers. For
sequence-to-sequence models, SHAP values can be computed at the token or
event level, allowing us to quantify how specific time steps influence
the prediction of future events. In transformer-based architectures,
SHAP can be applied to the \textbf{input embeddings}, intermediate
\textbf{hidden states}, or \textbf{attention scores} to reveal:
\end{quote}

\begin{itemize}
\item
  Which events or tokens had the greatest marginal impact on the
  predicted next event or outcome.
\item
  How attention mechanisms redistribute importance across the temporal
  context.
\item
  Whether the model's focus aligns with clinically relevant points in
  the treatment trajectory.
\end{itemize}

\begin{quote}
This token-level or event-level attribution is critical for medical
applications, where the interpretability requirement extends beyond
knowing \emph{which variables matter} to understanding \emph{when and in
what context they matter}.
\end{quote}

The \textbf{LIME} (Local Interpretable Model-agnostic Explanations)
framework {[}22{]} provides an intuitive method for interpreting
predictions made by complex machine learning models. Developed by
Ribeiro et al., \emph{LIME} is a model-agnostic approach that explains
individual predictions by approximating the black-box model with a
locally interpretable surrogate model, such as a linear regression.
LIME\textquotesingle s core philosophy is to make sense of "why" a model
produces a specific prediction by isolating the most influential
features for that decision.

\emph{LIME} works by perturbing the input data around the instance being
explained and observing the corresponding changes in the
model\textquotesingle s predictions. These perturbed samples are then
weighted based on their proximity to the original instance, and a
simpler, interpretable model (e.g., a linear or decision tree model) is
fitted to approximate the black-box model\textquotesingle s behavior in
this localized region. The weights assigned to each feature in the
surrogate model reveal their importance in driving the
model\textquotesingle s prediction.

\emph{LIME} is inherently flexible and can be applied to any machine
learning model, including linear models, tree-based models, neural
networks, and even ensemble methods. It is especially useful for models
that are inherently opaque, such as deep learning architectures or
boosted trees, where understanding feature contributions is challenging.
By generating explanations for individual instances, \emph{LIME} allows
practitioners to validate model behavior, identify biases, and ensure
alignment with domain knowledge.

When predicting the next event based on patient history, both
\emph{SHAP} and \emph{LIME} methods can highlight the key events or
intervals that significantly influence the prediction, enhancing
interpretability. We aim to leverage these insights to identify which
aspects of the protocol are most critical to treatment outcomes,
comparing results both with and without the new protocol intervention.
Ultimately, our goal is to refine the protocol based on practical
outcomes, integrating the assessed impact of actions taken or omitted
during treatment.

An intriguing example for the integration of such methods in medical
informatics is the article \emph{``Feature Importance to Explain
Multimodal Prediction Models: A Clinical Use Case''} {[}17{]}, which
developed a multimodal deep learning model to predict 30-day
postoperative mortality in elderly patients undergoing hip fracture
surgery. The model combined five data modalities, including
pre-operative static patient data (e.g., demographics, lab results, and
comorbidities), hip and chest X-rays, and per-operative data such as
vital signs and medications recorded during surgery. To process these
diverse data sources, ResNet50 was used for feature extraction from
X-rays, bi-directional LSTMs for sequential vital signs, and fully
connected neural networks for static patient data and medication
features. The extracted features were concatenated into a unified
representation, which was fed into a fully connected neural network to
predict mortality risk.

To enhance clinical interpretability, \emph{Shapley} values were
employed to explain the model's predictions. Global explanations showed
that pre-operative static data contributed most significantly,
accounting for about 70\% of the predictive power, while local
explanations provided individualized insights into the features that
influenced specific predictions. This approach demonstrates the
potential of integrating multimodal data with explainable AI methods to
support clinical decision-making while maintaining transparency in model
outputs.

Although their application of \emph{Shapley} values differs from our
intended use - since they focus on identifying important predictors for
outcomes, while we aim to take it a step further and to use these
insights on the temporal patterns to update \emph{GLs} - the goals
remain closely aligned - identifying key features that contribute to the
outcome prediction, creating actionable insights.

\section{3 Methods}\label{methods}

This chapter details the systematic approach used to develop predictive
models that incorporate temporal patterns and guideline adherence. The
methodology is structured into two main parts: \textbf{Computational
Methods} and \textbf{Evaluation Methods}.

\begin{itemize}
\item
  \textbf{Computational Methods} focuses on data preprocessing, model
  development, and training approaches.
\item
  \textbf{Evaluation Methods} describe how model performance is
  assessed, including comparison metrics and real-world validation.
\end{itemize}

By utilizing real-world data, the study aims to assess the effectiveness
of updated medical guidelines and their impact on patient outcomes.
Unlike other studies focused on medical outcome predictions, this
research targets short-term predictions for adult diabetic patients,
using a diverse dataset from six hospitals, each with unique medical
teams and practices. Furthermore, it specifically examines the role of
guideline adherence, as reflected by \emph{QA Scores}, and evaluates the
influence of new training initiatives on protocol implementation and
adherence in real-world hospital settings. Distinctively, this research
introduces a custom transformer-based architecture that has been
tailored to handle the temporal dynamics and clinical nuances of an
\emph{EMR}, addressing limitations in existing models and enabling
improved short-term predictions.

\includegraphics[width=6.26772in,height=3.04167in]{media/image17.png}

\textbf{Figure 3}: Methodology Overview - The figure illustrates the
research process, from knowledge base update to current protocol to the
comparison of different models and assessment of GLs importance to
treatment outcome.

\emph{TIRP = (Frequent) Time Interval Relation Patterns, TPF = Temporal
Probabilistic ProFiles, VSM = Vector Space Modeling, QATP = Quality
Assessment Temporal Patterns.}

\subsection{3.1 Computational Methods}\label{computational-methods}

The goal of this study is to predict hospitalization outcomes for
diabetic inpatients by leveraging temporal and sequential modeling
approaches. Our predictive models aim to forecast outcomes (e.g.,
discharge status, complications such as heart attack, stroke, or
mortality) at specific time points (up to the Nth hospitalization day)
based on patient data collected up to an earlier time point (Kth day),
where \textbf{48 hours ≤ 𝐾 \textless{} 𝑁 ≤ 14 days}.

This process involves a few major computational stages:

\begin{enumerate}
\def\labelenumi{\arabic{enumi}.}
\item
  \textbf{Temporal Data Preprocessing} -- Learning the OMOP DB model
  concept index, acquiring medical knowledge (KBTA) and processing the
  model into a concept based temporal interval dataset.
\item
  \textbf{Temporal Data Abstraction} -- Extracting structured,
  meaningful representations from raw EMR data including events, states
  and trends.
\item
  \textbf{QA Scores Calculation} -- Calculating adherence to clinical
  guidelines scores using fuzzy logic, based on the \emph{Quality
  Assessment Temporal Patterns (QATPs)}.
\item
  \textbf{Predictive Model Development} -- Building
  temporal-frequent-patterns-based and transformer-decoder models using
  machine learning and deep learning techniques.
\item
  \textbf{Optimization} -- Fine-tuning models for optimal performance
  including hyperparameters and loss optimization, augmentation,
  variance insertion and feature selection.
\end{enumerate}

The algorithms and systems discussed in this chapter were previously
covered in the \emph{Related Work} chapter \textbf{{[}Chapter 2{]}}, so
this section focuses solely on their intended application within this
research.

\subsubsection{3.1.1 Temporal Data Abstraction and QA Scores
Calculation}\label{temporal-data-abstraction-and-qa-scores-calculation}

A crucial aspect of this study is the development of a structured
temporal abstraction framework, which transforms raw electronic medical
records (EMR) into clinically meaningful states. This process is
knowledge-driven, leveraging a formally defined \emph{Knowledge Base
(KB)} that encodes diabetes treatment guidelines. Key stages:

\textbf{Building Database Specific Index}

\begin{itemize}
\item
  Identifying relevant concepts for measurements, medications and
  observations and building a dataset specific index, as the
  \emph{Kineret} Dataset {[}34{]} is not mapped to any international
  coding system.
\item
  Automating the preprocess from \emph{OMOP} model to temporal dataset
  input, using the hand crafted index.
\item
  Knowledge acquisition regarding discretization thresholds, event
  relevance durations and definitions for \emph{KBTA} abstractions using
  medical knowledge and clustering methods (K-MEANS, with KDE
  distribution plot) to identify density peaks and dips in every
  feature. The final discretization thresholds are decided using the
  discovered thresholds combined with the medical knowledge to create
  reasonable abstractions.
\end{itemize}

\textbf{Knowledge Base (KB) Development and Acquisition}

\begin{itemize}
\item
  We construct a machine-readable \emph{KB} that encodes key elements of
  diabetes treatment guidelines (\emph{GLs}), specifically tailored to
  the diabetic population (ages 18+) in the \emph{Kineret Dataset.} This
  \emph{KB} includes declarative definitions such as ``high blood
  pressure'' or ``prolonged hyperglycemia,'' allowing automated
  abstraction of raw medical data into discrete states and written in
  TAK language. A developed helper automation for this process is
  described in \textbf{{[}Appendix 2{]}}.
\item
  The diabetes \emph{GLs} used in this study are part of a new national
  protocol for hospital-based diabetes care.
\end{itemize}

\textbf{Guideline Compliance Encoding}

\begin{itemize}
\item
  A \emph{Fuzzy Logic} grading system is applied using QATPs to assess
  the level of adherence to medical guidelines for each abstracted
  pattern during treatment. For example, the existence of a recommended
  insulin dosage up to t hours after a meal can be enforced using
  pattern abstractions and aggregated as ``insulin adherence score''
  using \emph{Fuzzy Logic}.
\item
  Unlike real-time decision-support systems (e.g., Mobiguide {[}7{]}),
  our approach does not aim to update procedural knowledge or provide
  immediate recommendations but rather enables light retrospective
  scoring and assessment, that is integrated to assess the improvement
  in adherence but mainly its effects on the outcomes.
\end{itemize}

\textbf{Temporal Abstraction and Data Transformation}

\begin{itemize}
\item
  We utilize the Temporal Mediator to process raw, time-dependent data
  into structured time-interval-based states.
\item
  Each patient's \emph{EMR} is segmented into meaningful intervals,
  reducing noise and preserving clinically relevant trends while
  avoiding unnecessary dimensionality.
\item
  The output of this stage consists of a sequence of discrete
  abstractions (analogous to words in text processing), where some
  abstractions (patterns) are associated with a fuzzy compliance score
  based on adherence to guidelines.
\end{itemize}

\subsubsection{3.1.2 Predictive Model
Development}\label{predictive-model-development}

To predict short-term medical outcomes, we explore two primary modeling
approaches: \textbf{Frequency-Based Temporal-Patterns Vector-Space
Modeling (TPF/TIRP+VSM)} which is based on the horizontal support and
distribution of \emph{Time Interval Relation Patterns (TIRPs} and
TPFs\emph{)} in a patient's records, and \textbf{Transformer Modeling}
which uses the \emph{KBTA} output as tokenized input to predict an
expected events trajectory. Each method requires different data
representations and preprocessing, which are detailed below.

The predictive models will utilize the diverse set of patient data from
the \emph{Kineret Dataset}, including:

\begin{itemize}
\item
  Patient background (e.g., age, medical history)
\item
  Clinical measurements and lab results
\item
  Administered medications and dosages
\item
  Observations made during hospitalization
\item
  Quality-assessment scores derived from staff adherence to clinical
  guidelines
\end{itemize}

\paragraph{3.1.2.1 Frequency-Based Temporal-Patterns Vector-Space
Modeling
(TPF/TIRP+VSM)}\label{frequency-based-temporal-patterns-vector-space-modeling-tpftirpvsm}

In this approach, abstracted patient states and compliance scores are
transformed into structured feature vectors suitable for classical
machine learning models. All three methods described below use
KarmaLego's TIRP output, computed on the time-abstracted patient states
generated in Section 3.1.1. These TIRPs represent frequent sequences of
clinical states (e.g., ``prolonged hyperglycemia'' → ``insulin dose
increase'' → ``hypoglycemia'').

We will construct three alternative vector-space representations:

\begin{enumerate}
\def\labelenumi{\arabic{enumi}.}
\item
  \textbf{Basic TIRP Cardinality Vector} -- For each patient at the
  \emph{K}th day, count the \textbf{horizontal support} (number of
  occurrences) of each frequent TIRP in their sequence. This assumes
  that the hospital stay exceeds \emph{K} days for all patients.
\item
  \textbf{Relative-Support TPF} -- For each patient, compute the
  horizontal support of each frequent TIRP, then normalize it relative
  to the same TIRP's support in the rest of the cohort using min--max
  scaling to the range {[}0, 1{]}. This captures whether a pattern is
  unusually prominent for a specific patient.
\item
  \textbf{Duration-Normalized TPF} -- For each frequent TIRP, scale its
  contribution in the feature vector by the proportion of the total
  input duration (\emph{K} days) that the pattern spans. This emphasizes
  patterns that occupy longer portions of the patient's treatment
  trajectory.
\end{enumerate}

These summary vectors capture a representation of the visit's context
and can be fed into models such as K-Nearest Neighbors (KNN), XGBoost,
or Deep Neural Networks (DNNs). Additionally, guideline compliance
scores will be incorporated into these representations, by using the
distribution of ``correct'' or ``incorrect'' treatment patterns in these
vectors. Feature selection methods can guide us on which patterns are
more predictive than others. All vector-space approaches will be
compared to determine which best captures meaningful trends in the
short-term prediction of diabetes treatment outcomes.

\paragraph{3.1.2.2 Sequential Models}\label{sequential-models}

Unlike TPF/TIRP+VSM approaches, sequential models retain the structured
state representations and operate directly on the \textbf{temporal
sequence of states trends and patterns} rather than transforming them
into fixed feature vectors of pattern distribution. This method is
motivated by the need to capture \textbf{not only if but also when} an
event is likely to occur, as well as to output a sequence of outcomes,
if applicable.

\begin{itemize}
\item
  \textbf{Next-State Prediction for Temporal Sequence Modeling}

  \begin{itemize}
  \item
    Each state/event sequence will be treated analogously to a word
    sequence in NLP, where models attempt to predict the next expected
    state based on past patterns.
  \item
    An additional Δt prediction head is added to the model, to assess
    the timing of the next event. This Δt is crucial for the transformer
    generative action's positional encoding, but in the evaluation
    events timing might be assessed using the MEAL events rather than
    the attempted Δt, depending on the actual learning quality of this
    head.
  \item
    The goal is to develop a sequential prediction model that learns
    expected event trajectories and timings, distinguishing normal
    progressions from those leading to complications. At every step (t),
    the model is trained to predict the next most probable
    (deterministic) event-token, while respecting constraints on illegal
    tokens at a specific point in time (For example, glucose cannot be
    \ul{high} and \ul{medium} at the same time).
  \item
    Unlike classification-based approaches, this method can include the
    timing of outcome predictions rather than merely predicting their
    occurrence, and is inherently multi-label approach, allowing the
    same model to predict multiple outcomes for the same hospital visit.
  \end{itemize}
\item
  \textbf{Data Preprocessing for Sequential Models}

  \begin{itemize}
  \item
    This approach requires a structured representation of a
    patient\textquotesingle s temporal intervals. The temporal input
    dataset for this section (after processed by the Mediator) will have
    a structure like this example:
  \end{itemize}
\end{itemize}

\begin{quote}
\includegraphics[width=4.49479in,height=0.72967in]{media/image29.png}
\end{quote}

\begin{itemize}
\item
  Such a dataset will include multiple dimensions (for every feature)
  which cannot fit into a singular input sequential model. To handle the
  dimensionality while still allowing intervals to have a duration (in
  the input) every interval with duration (states / trends / patterns)
  is divided to 2 tokens (``\_Start'', ``\_End''). For example,
  GLUCOSE\_STATE\_High will be tokenized as GLUCOSE\_STATE\_High\_Start,
  \ldots, GLUCOSE\_STATE\_High\_End, depending on the relative timing of
  the interval closing point to other intervals. The correctness of the
  abstraction process guarantees that no sequence will be corrupted
  (e.g., GLUCOSE\_STATE\_High\_Start, \ldots,
  GLUCOSE\_STATE\_High\_Start, \ldots, GLUCOSE\_STATE\_High\_End
  resulting in a closing token for 2 identical starting tokens).
\item
  In addition, in order to derive dependencies in time, for every token
  we'll calculate the Δt from the Admission token (global dependency).
  This makes the current shape of the dataset to be like in this
  example:
\end{itemize}

\begin{quote}
\includegraphics[width=3.30729in,height=0.92366in]{media/image7.png}
\end{quote}

\begin{itemize}
\item
  While solving the interval duration representation problem, this
  method makes 2 tokens of the same concept to be treated differently,
  which has the potential to damage the informative qualities we aim to
  achieve from the embedding. To handle that, 4 levels of tokenization
  were created for every event token:

  \begin{itemize}
  \item
    Raw Concept (e.g., GLUCOSE)
  \item
    Concept / TAK (e.g., GLUCOSE\_STATE)
  \item
    Concept + Value (e.g., GLUCOSE\_STATE\_HIGH)
  \item
    Concept + Value + Position (e.g., GLUCOSE\_STATE\_High\_START)
  \end{itemize}
\end{itemize}

\begin{quote}
This method allows information to be carried across different tokens of
the same event, and even different events of the same raw-concept. For
example, the transformer will now be informed that
GLUCOSE\_STATE\_High\_END is the closing token for previously seen
GLUCOSE\_STATE\_High\_START and that GLUCOSE\_TREND shares at least some
(learnt) similarity to GLUCOSE\_STATE. A tokenizer (EMRTokenizer) is
created to handle the encoding and decoding of the new medical-events
``language''.
\end{quote}

\begin{itemize}
\item
  Finally, a demographic dataset ({[}CTX{]}), holding Age, Gender,
  number of prior hospitalizations in 6 and 12 months, and past
  occurrences (before current admission) of the relevant medical
  complications is created, and scaled using a standard scaler (fitted
  on train, used on test). This vector will be accompanied by the
  \emph{Charlson Comorbidity Index (CCI)} that quantifies a patient's
  overall disease burden by assigning points to comorbid conditions
  based on their severity and mortality risk, producing a single score
  that reflects how ill the patient is. This score is considered highly
  informative.
\item
  We aim to compare 3 inputs for the same architecture:
\end{itemize}

\begin{enumerate}
\def\labelenumi{\arabic{enumi}.}
\item
  \textbf{A base-line} of token representation is curated by
  equal-frequency discretization of the raw input data, which means
  these tokens will not be medical knowledge based and \textbf{will not
  have a duration}.
\item
  \textbf{An Interval based} token representation curated with the
  described Mediator's temporal abstractions using only events, states
  and trends (not incorporating treatment adherence patterns).
\item
  \textbf{Interval + Patterns based} token representation which will
  encompass both the abstracted EMR and treatment adherence patterns
  (\emph{QATP}s) found in the input, which can inform the model of what
  patterns should receive more attention weight.
\end{enumerate}

\begin{itemize}
\item
  \textbf{Embedding method}

  \begin{itemize}
  \item
    Each event is decomposed into four hierarchical levels of token
    information as described above. These four components are embedded
    separately using dedicated embedding layers, each mapping the tokens
    into the same fixed-size embedding space (defined by the final
    embedding dimension, e.g., 128).
  \item
    To incorporate temporal patterns, \textbf{Time2Vec} layer is used
    for \textbf{absolute timing} (time from admission). This encoder
    outputs a vector containing a linear component and multiple periodic
    components (sine functions) as described in the paper {[}19{]}.
    These embeddings are concatenated and then projected to the same
    final embedding dimension via a linear layer, ensuring that temporal
    information is harmonized with the token-level embeddings.
  \item
    The final event representation is created by \textbf{concatenating}
    the four token embeddings (raw concept, concept type, concept-value,
    concept-value-position) with the projected time embedding. This
    results in a combined vector of \textbf{5 × embedding dimension}
    (e.g., 5 × 128 = 640). A final \textbf{linear projection layer} then
    compresses this vector into the predefined embedding dimension
    (e.g., 128), ensuring a consistent, compact representation for every
    event.
  \item
    In addition to event embeddings, a dedicated \textbf{{[}CTX{]}
    token} (learnable) is prepended to every patient's event sequence.
    This token is combined with a projection of the patient's static
    context data (age, gender, etc.) and possibly the adherence QA
    aggregated score to encode global patient characteristics that
    modulate the interpretation of dynamic events.
  \item
    During \,training \textbf{phase\,1}
    the\,\emph{EMREmbedding}\,network is trained with three
    self‑supervised losses that act in concert. First, Pytoch's
    \textbf{Multi‑Label Binary Cross Entropy (BCE) with logits} {[}31{]}
    objective predicts the \emph{set} of events likely to appear within
    a short look‑ahead window, so the model learns which events can
    legitimately follow the current one while remaining tolerant to
    slight swaps in order. Second, a \textbf{masked‑language‑modeling
    (MLM)} {[}30{]} task, applied with the classic
    80\,\%\,→\,{[}MASK{]}, 10\,\%\,→\,random, 10\,\%\,→\,keep strategy
    as in the original BERT paper, forces the encoder to reconstruct
    hidden tokens from their bidirectional context, capturing
    longer‑range dependencies that a strict next‑token loss would miss.
    Finally, an auxiliary \textbf{Mean Square Error (MSE)} {[}31{]} loss
    on a time‑prediction head teaches the model to estimate when the
    next event should occur, giving it temporal intuition. Together,
    these three signals encourage the embedder to internalise the core
    clinical co‑occurrence patterns, their approximate timing, and the
    broader semantic context, all while being robust to the inevitable
    noise and minor ordering irregularities present in real‑world
    electronic medical records.
  \end{itemize}
\item
  \textbf{Transformer Method}

  \begin{itemize}
  \item
    Foundation \textbf{in GPT-2 Architecture:\\
    } The transformer model is built upon the GPT-2 architecture by
    Radford et al. (2019) {[}28{]} and further refined in the
    \textbf{nanoGPT} implementation by Karpathy {[}29{]}. It uses
    stacked self-attention and feed-forward layers with residual
    connections to model event sequences autoregressively.
  \item
    \textbf{Replacement of Token/Positional Embeddings:\\
    } Unlike standard GPT-2, our transformer operates on the
    \textbf{context-aware event embeddings} produced by the
    \textbf{EMREmbedding} module. These embeddings already incorporate
    temporal, semantic, and patient-specific information, removing the
    need for separate token and positional embeddings. The model
    receives a trained embedder object as input with the possibility to
    tune its weights during training.
  \item
    \textbf{Model Stack and Residual Scaling:\\
    } The transformer stack consists of multiple causal self-attention
    blocks and gated MLPs, each with layer normalization. A residual
    scaling factor of
    \includegraphics[width=0.61979in,height=0.37537in]{media/image11.png}
    is applied to each residual branch to ensure stable gradient flow
    during deep model training.
  \item
    \textbf{Dual Prediction Heads:}

    \begin{itemize}
    \item
      The \textbf{next-token prediction head} uses weight tying with the
      event embedding layer, maintaining parameter efficiency.
    \item
      An additional \textbf{Δt prediction head} is trained to predict
      the timing of the next event (Δt), allowing explicit modeling of
      temporal intervals alongside event predictions.
    \end{itemize}
  \item
    \textbf{Customized Multi-Hot Focal BCE Loss:\\
    } The training objective uses a \textbf{Multi-Hot Focal BCE Logits
    Loss} formulation. Instead of strict next-token prediction, this
    method predicts a \textbf{set of relevant events} within a future
    window, while balancing rare tokens in the loss computation,
    reflecting the real-world noise and flexibility in clinical event
    timing. The method takes into account only future events within
    window k that \textbf{are legal} in the current context up to event
    t (including t). Meaning, the model does not propagate loss over
    tokens that are not relevant to the context {[}0 → t{]}. This is a
    method based on complete teacher's forcing where at every step the
    model receives as context only ground truth tokens, and as targets
    only tokens that does not violate structural validity to choose
    from.
  \item
    \textbf{Incorporation of Clinical Penalties:\\
    } The transformer training loss includes \textbf{domain-specific
    penalty functions} to enforce clinically meaningful predictions:

    \begin{itemize}
    \item
      \textbf{Meal order penalty:} Penalizes predictions that violate
      daily meal cycles (Breakfast → Lunch → Dinner → Night (snack) →
      Breakfast). Note that the dataset does not offer the received
      meals information, so a fixed schedule (7:00, 12:00, 18:00, 22:00)
      is assumed during data preprocessing.
    \item
      \textbf{Hallucinated interval penalty:} Penalizes illegalities in
      the shape of the generated intervals. Terms:
    \end{itemize}
  \end{itemize}
\end{itemize}

\begin{quote}
If token== \textquotesingle GLUCOSE\_TREND\_inc\_START\textquotesingle,
base(tok) == \textquotesingle GLUCOSE\_TREND\_inc\textquotesingle,
concept(tok) == \textquotesingle GLUCOSE\_TREND\textquotesingle{}
\end{quote}

\begin{itemize}
\item
  END is illegal if the base is not open yet (You can\textquotesingle t
  see \textquotesingle GLUCOSE\_TREND\_inc\_END\textquotesingle{} before
  opening it
  \textquotesingle GLUCOSE\_TREND\_inc\_START\textquotesingle). Enforced
  using base(tok).
\item
  START is illegal if the base already open (You can\textquotesingle t
  have \textquotesingle GLUCOSE\_TREND\_inc\_START\textquotesingle{}
  after \textquotesingle GLUCOSE\_TREND\_inc\_START\textquotesingle{}
  without seeing
  \textquotesingle GLUCOSE\_TREND\_inc\_END\textquotesingle). Enforced
  using base(tok).
\item
  START of concept(tok) is illegal if concept(tok) is still open
  (meaning you can\textquotesingle t have
  \textquotesingle GLUCOSE\_TREND\_inc\_START\textquotesingle{} after
  \textquotesingle GLUCOSE\_TREND\_dec\_START\textquotesingle{} without
  seeing \textquotesingle GLUCOSE\_TREND\_inc\_END\textquotesingle).
  Enforced using concept(tok).
\end{itemize}

\begin{itemize}
\item
  \textbf{Total Loss Composition:}
\end{itemize}

\begin{quote}
The final training loss is a weighted sum of the next-token \textbf{BCE}
loss, the \textbf{structured penalties} (meal order, interval
correctness), and the \textbf{Δt MSE} regression loss + a minor penalty
to encourage time monotonicity. Each component is weighted according to
its clinical relevance and sparsity, ensuring the model learns not only
accurate event predictions but also clinically plausible event sequences
with precise temporal modeling. This training process is referred to as
\textbf{phase 2}.
\end{quote}

\begin{itemize}
\item
  The model defines a \textbf{warm up} phase in the training (e.g., 10
  epochs) where no stopping criteria can be triggered and a
  \textbf{patience} phase (e.g., 5 epochs) to truly decide if a model is
  not improving. During this warmup phase the weight of the penalties
  increases gradually to allow the model to first learn the vocab and
  patterns, and only then to slowly introduce time and structure.
\item
  The model also \textbf{masks} non-structural tokens (not ADMISSION,
  MEALs, PAD etc) randomly in a growing ratio on the input to encourage
  robustness. This method is called Curriculum by Masking (CBM).
\item
  In addition, the data loader oversamples admissions with rare outcomes
  (e.g., DEATH, \textless5\%). Combined with the weighted loss and the
  CBM, this method injects variance while balancing the influence of
  rare medical events, which are common in medicine.
\end{itemize}

We call this model \emph{Transform-EMR} {[}32{]}, and its current
(updating) implementation is available on GitHub. Currently, the planned
architecture is only \textbf{Transformer-based}, due to its appealing
attention mechanism for this task and its training parallelism over the
block, but depending on our progress, results and time constraints, we
might also attempt an \textbf{LSTM} architecture with similar loss and
targets for benchmarking.

\paragraph{3.1.2.3 Comparison and Research
Goals}\label{comparison-and-research-goals}

By comparing these approaches, we aim to determine:

\begin{itemize}
\item
  Whether \textbf{TPFs combined with guideline compliance scores} as
  baseline for vector space models (VSMs) provide a superior predictive
  signal for short-term treatment outcomes. These models benefit from
  clinically pre-defined feature sets, potentially leveraging
  well-established patterns of care.
\item
  Whether our \textbf{transformer architecture} can outperform VSMs in
  both \textbf{accuracy and timing of predictions}, particularly in
  anticipating when an adverse event is likely to occur. To assess the
  influence of guideline adherence, the transformer will be trained
  \textbf{with and without} the compliance scores. Additionally, we will
  explore \textbf{tokenizing guideline patterns directly} within the
  transformer input, potentially guiding the model to learn more robust
  clinical causality.
\item
  Unlike TPFs/TIRPs, which rely on predefined/ explicitly discovered
  patterns, the transformer has the capacity to \textbf{discover novel,
  undefined patterns} that may influence patient survival, thanks to its
  ability to model complex, hierarchical dependencies through attention
  mechanisms.
\end{itemize}

The results will inform which temporal modeling strategy is most
effective for real-world deployment in clinical decision support.

\subsubsection{3.1.3 Hyperparameter
Optimization}\label{hyperparameter-optimization}

Optimizing model performance is crucial for ensuring that we maximize a
predictive algorithm's potential, by adjusting the way it learns the
data representation. To achieve this, we will employ \textbf{automated
hyperparameter optimization} using Optuna {[}21{]}, a powerful framework
for efficiently searching hyperparameter spaces, and a qualitative
search over diverse predefined hyperparameter space, in places where
Optuna is infeasible (Deep Sequential Models).

\textbf{Optimization Strategy}

Each of the modeling approaches outlined in \textbf{{[}Chapters
3.1.2{]}} will require fine-tuning of key hyperparameters to achieve the
best predictive performance. The optimization process will include:

\textbf{Vector-Space Machine Learning Models}

\begin{itemize}
\item
  \textbf{KNN:} Optimization of k-value selection, distance metrics
  (e.g., Euclidean, Cosine Similarity, iDTW), and weighting functions.
\item
  \textbf{XGBoost:} Fine-tuning of tree depth, learning rate, number of
  estimators, and feature importance thresholds.
\end{itemize}

\textbf{Sequential Models}

\begin{itemize}
\item
  \textbf{LSTMs and Transformers:} Optimization of hidden state size,
  number of layers, attention heads, learning rate schedules, masking
  and loss function to balance generalization and efficiency.
\item
  \textbf{Training Parameters:} Batch size, block size, learning rate,
  weight decay, dropout, number of epochs, BCE window (number of tokens
  per prediction) and penalties weights.
\end{itemize}

This optimization framework will ensure that each model operates at peak
efficiency, maximizing predictive accuracy across different
methodologies. By leveraging adaptive search strategies instead of
exhaustive grid search, we aim to reduce training time while maintaining
optimal model performance.

\subsection{3.2 Evaluation Methods}\label{evaluation-methods}

This study aims to address the following \textbf{key questions}:

\begin{itemize}
\item
  Can Transformers or other sequential models outperform interval-based
  temporal frequent pattern methods in medical short-term outcome
  predictions, as outlined in this study? Can \emph{KBTA} / \emph{TIRP}
  based tokenization help?
\item
  How does adherence to medical guidelines influence the outcomes during
  hospital stays?
\item
  Which components of the standardized protocols are most effective, and
  where do gaps or inconsistencies in their implementation arise?
\item
  How guideline-compliance scores have been influenced by recent
  interventions aimed at improving clinical care.
\end{itemize}

The \textbf{hypotheses} underpinning this study are:

\begin{itemize}
\item
  Patients treated with higher adherence to standardized diabetes
  protocols will experience lower short‑term in‑hospital complication
  rates than patients with lower adherence.
\item
  A Transformer‑based model will outperform a frequent interval-based
  temporal patterns vector-space baseline model in predicting short‑term
  complications, indicating that it captures clinically relevant
  patterns not encoded in handcrafted features.
\item
  Knowledge-based temporal abstraction before tokenization (states/
  trends/ patterns) will yield higher predictive performance than
  discrete equal-frequency tokenization.
\item
  Model‑derived feature‑importance ranking will highlight specific
  protocol components (e.g., timely insulin titration, glucose
  monitoring frequency) as dominant contributors to reduced
  complications and will also highlight protocol components that suffer
  from implementation gaps.
\item
  Mean guideline‑compliance scores will increase following the recent
  clinical interventions.
\end{itemize}

The study's objectives are as follows:

\begin{enumerate}
\def\labelenumi{\arabic{enumi}.}
\item
  \textbf{Knowledge Base Update and Temporal Abstraction:} Develop a
  comprehensive knowledge base to support the implementation of the
  updated diabetes care protocol in future studies, as well as to
  support the QA scores calculation for this study using the existing
  temporal abstraction algorithms.
\item
  \textbf{Medical Outcome Prediction:} Utilize the discrete and
  abstracted data to predict hospitalization outcomes at specific time
  points (up to Nth hospitalization day), based on prior data collected
  up to the Kth day (48 hours ≤ K \textless{} N ≤ 14 days). This
  prediction task will be conducted using the specified TIRP/ TPF based
  methods (Section 3.1.2.1) and the transformer based methods (Section
  3.1.2.2).
\item
  \textbf{GL Adherence as Features:} Introduce guideline adherence
  scores as predictive features and evaluate their influence on the
  medical outcomes.
\item
  \textbf{Protocol Adjustment:} Evaluate the effectiveness of care
  guidelines, identifying under-implemented areas and components with
  the strongest connection to specific outcomes. Highlight opportunities
  for refinement and improve consistency in implementation.
\item
  \textbf{Impact of Interventions:} Assess the effect of two
  interventions across six hospitals - {[}1{]} the introduction of a
  diabetes GL learning software module for medical teams, and {[}2{]}
  the introduction of a diabetes quality measure for managing
  hyperglycemia events - on the compliance scores of the medical staff.
  This will be achieved by comparing \emph{GL}-based \emph{QA} scores
  before and after these interventions.
\end{enumerate}

The dataset, referred to as the ``\emph{Kineret dataset}'' {[}34{]} is
curated of 6 Israeli government hospitals data: Shamir, Barzilai, Bnei
Zion, Tzafon (Poria), Galilee (GMC), Hillel Yaffe (HY). The study cohort
is defined by the following inclusion and exclusion criteria:

\begin{itemize}
\item
  \textbf{Age:} Inpatients aged 18 and above.
\item
  \textbf{Department Exclusions:} Excludes patients from obstetrics and
  pediatrics departments.
\item
  \textbf{Diabetes Indicators:} Includes patients diagnosed with
  diabetes or with elevated blood glucose levels (≥180 mg/dL twice or
  ≥250 mg/dL once during the same hospitalization).
\item
  \textbf{Hospital Stay Duration:} Hospitalizations lasting 48 hours to
  14 days.
\end{itemize}

The dataset includes patient background information, clinical
measurements, administered medications, observations, and outcomes such
as discharge status or complications (e.g., heart attack, stroke, or
death). Two datasets were collected:

\begin{enumerate}
\def\labelenumi{\arabic{enumi}.}
\item
  June 2022 to August 2023, representing the pre-protocol implementation
  phase.
\item
  August 2023 to January 2025, reflecting data after the protocol was
  updated.
\end{enumerate}

This dataset is hosted externally on a secure, cloud-based platform to
ensure data privacy and compliance with the Ministry of Health
regulations. Deploying the required algorithms - such as the
\emph{Temporal Mediator}, \emph{KarmaLego} and the \emph{Custom
Transformer} - presented several engineering challenges. We needed to
adapt these systems for compatibility with Linux-based containers and
optimize them for deployment in resource-constrained environments.
Additionally, building robust data extraction pipelines within the
remote environment, powered by our custom concept codes index that had
to be derived from the actual data was crucial for ensuring the seamless
flow of clinical data into these containers while managing computational
and memory limitations. These efforts were essential for enabling
temporal abstraction and pattern mining without compromising performance
or reliability.

\subsubsection{3.2.1 Evaluation of the
Knowledge-Base}\label{evaluation-of-the-knowledge-base}

As described in {[}\textbf{Section 3.1.1}{]}, The first steps of the
project are focused on the acquisition informative abstractions and the
medical guidelines in machine readable manner, to be fed into the
Temporal Mediator, creating the temporal abstractions which are then
used as features and for the calculation of the \emph{Quality Assessment
Score} per action based on the \emph{Fuzzy Logic}. Given that creating a
treatment quality assessment system (like in the \emph{Picard} project)
is not a part of this study's goals the Quality-Assessment Knowledge
implementation (\textbf{Task No.1 in Section 3.2}) will be evaluated
heuristically, and Technical / Functional Evaluations {[}2{]} will not
be needed.

As described, when creating the abstractions parameters, a few methods
were combined with public medical knowledge (Internet + medical
professional assistance):

\begin{enumerate}
\def\labelenumi{\arabic{enumi}.}
\item
  Calculation of 1D K-Means for every concept (e.g., administered Bolus
  Insulin Dosage) in order to identify 2-5 meaningful clusters, taking
  the middle point between their centroids as discretization thresholds.
\item
  Visualization using a KDE density histogram, showing peaks and dips in
  administered dosage frequency in the cohort.
\item
  Qualitatively combining these found thresholds with medical knowledge
  to decide on the final threshold for every concept.
\end{enumerate}

\includegraphics[width=6.26772in,height=2.625in]{media/image19.png}

\textbf{Figure 4:} An example of Bolus Insulin density plot, showing 3
main density peaks. The red lines are dips markers (local minimum) and
the purple lines represent the final thresholds found for this concept.

\begin{enumerate}
\def\labelenumi{\arabic{enumi}.}
\setcounter{enumi}{3}
\item
  In addition to the state discretization thresholds, relevance
  durations were also needed (e.g., the duration after admission a
  medication is considered in the blood stream) and trend limits (what
  is considered a medically significant increase/decrease in a dosage or
  measurement). A qualitative literature search was done in order to
  decide on these.
\end{enumerate}

After these rules were decided, the resulting temporal abstractions were
also evaluated but in an empirical manner, mostly to validate that they
were created properly relative to the real data distribution. Meaning,
the medical knowledge can instruct us on the clinically approved normal
ranges for measurements, but if their real world ranges vary
significantly from these recommendations (for example, high density on
out-of-range values or ``starved'' discrete values) - adjustments need
to be made. These were evaluated after abstracting the data, and
included:

\begin{itemize}
\item
  Checking the avg. duration of each concept state/ trend.
\item
  Checking the ratio of temporal abstractions with duration compared to
  events (to confirm concepts were abstracted properly).
\item
  Checking that every pre-defined discrete state is represented after
  abstraction and was not excluded due to wrong pre-defined
  discretization criterias.
\end{itemize}

As these evaluations are not a main part of this study and mainly used
for internal validation, they are shown in \textbf{{[}Appendix 1{]}}.

\subsubsection{3.2.2 Evaluation of the Predictive
Models}\label{evaluation-of-the-predictive-models}

Building on the predictive modeling approaches outlined in
{[}\textbf{Section 3.1.2{]}}, this section focuses on assessing model
performance to ensure reliability and clinical applicability. The
evaluation process is designed to validate the predictive ability of the
models while considering the complexity and evolving nature of the
prediction task. Our full prediction targets list is:

\begin{itemize}
\item
  Release
\item
  Death (up to 30 days from release).
\item
  Complication: Ketoacidosis, Kidney disorder, Coma, Eye disorder,
  Hypoglycemia or Hyperglycemia (mild / severe), Cardiovascular
  disorder, Infection, Neurovascular complication.
\item
  As a byproduct of our method we'll also try to predict the visit's
  total duration.
\end{itemize}

{\def\LTcaptype{none} % do not increment counter
\begin{longtable}[]{@{}
  >{\raggedright\arraybackslash}p{(\linewidth - 2\tabcolsep) * \real{0.5936}}
  >{\raggedright\arraybackslash}p{(\linewidth - 2\tabcolsep) * \real{0.3193}}@{}}
\toprule\noalign{}
\endhead
\bottomrule\noalign{}
\endlastfoot
\textbf{Target} & \textbf{Vertical Support Cohort} \\
Release & 96.22\% \\
Death & 3.84\% \\
Hyperglycemia (glucose \textgreater{} 250 mg/dL) & 69.60\% \\
Severe Hypoglycemia (glucose \textless= 54 mg/dL) & 5.64\% \\
Ketoacidosis & 2.54\% \\
Kidney disorder & 12.69\% \\
Coma & 1.04\% \\
Eye disorder & 1.94\% \\
Cardiovascular disorder & 10.99\% \\
Infection & 5.77\% \\
Neurovascular complication & 5.45\% \\
\end{longtable}
}

\textbf{Table 1}: The distribution of the prediction targets within the
cohort. While terminal targets (release / death) can only appear once,
complications can appear multiple times in a single hospital visit. The
ratio shows the number of patients with at least 1 occurrence (vertical
support) of this target in their full visit trajectory.

The models are programmed to predict only one terminal outcome per
hospital visit (release/ death) but can predict multiple complications
to happen between the admission and that terminal outcome. There are 3
evaluation levels we employed:

\begin{itemize}
\item
  Were we able to predict \textbf{\ul{a}} future complication and
  outcome? (Easy)
\item
  Were we able to predict \textbf{\ul{the}} future complications and
  outcome? (Medium)
\item
  Were we able to predict \textbf{\ul{the}} future complications and
  outcome \textbf{\ul{at the}} correct time? (Hard, only possible for
  the generative model)
\end{itemize}

In the temporal frequency vector-space models prediction is
straight-forward (binary multi-label vector as output). For the
generative model we will create such a vector. Iterating over the
generated events, we will collect all outcomes and terminal outcomes
generated by the model (per patient) and will mark them in a binary
multi-label vector, depending on them meeting the criteria (for example,
max deviation of 24 hours from the real event time).

To ensure the reliability and accuracy of the models (\textbf{Task No.2
in Section 3.2}) and to assess the impact of adherence scores on
predictive performance (\textbf{Task No.3 in Section 3.2}), each model
will be evaluated using a range of performance metrics and validation
methods. These evaluations will determine the effectiveness and
generalizability of the models while also providing insights into how
guideline adherence contributes to prediction quality.

\begin{itemize}
\item
  \textbf{Precision}: Measures the proportion of true positive
  predictions among all positive predictions made by the model.
\item
  \textbf{Recall / Sensitivity}: Evaluates the model\textquotesingle s
  ability to identify all relevant instances within the data.
  Specifically, it measures the proportion of actual positive instances
  that the model correctly identifies as positive.
\item
  \textbf{F1 Score}: The harmonic mean of precision and recall,
  providing a balanced measure of both. Particularly useful with
  imbalanced datasets, as it ensures that both false positives and false
  negatives are considered.
\item
  \textbf{Specificity}: Measures the proportion of true negative
  instances correctly identified by the model. It reflects the model's
  ability to avoid false positives and is especially important when the
  cost of a false alarm is high.
\end{itemize}

These metrics will be calculated for micro and macro averages. In the
micro average, individual class contributions are weighted by their
frequency, treating each prediction as equally important. The macro
average, on the other hand, calculates the metric independently for each
class and then averages the results, treating each class equally
regardless of size. This differentiation is irrelevant for the
\emph{Accuracy} measurement.

The models will be evaluated using a \textbf{holdout test set}
consisting of 1000 patients. The test set is relatively small due to
computational resources restrictions in the sequential token-by-token
inference of a transformer model. This approach is chosen because deep
learning models are typically \textbf{less suited for K-fold
cross-validation}, a technique that partitions the data into multiple
subsets (folds) for repeated evaluation. Instead, a single test set
provides a practical and efficient way to assess model performance on
previously unseen data, while avoiding the computational burden of
multiple training cycles inherent in K-fold cross-validation.

In addition, comparison tests will be conducted to explore the impact of
different model and data configurations as described in \textbf{Sections
3.1.2.1} and \textbf{3.1.2.2}. This will help assess whether the
inclusion of QA data - reflecting adherence to updated guidelines -
improves the model's predictive power, and which modeling technique
works best for this prediction task.

Finally, the study will benchmark its results against existing research
in the field of medical outcome prediction for diabetes treatment. By
comparing our findings with those of similar studies, we aim to validate
the effectiveness of our model and understand how it performs relative
to established approaches. However, we acknowledge that our dataset is
not publicly available (unlike datasets such as MIMIC), which may limit
the comparability of our results and introduce variability due to the
dataset's uncontrolled quality. If needed to further validate our novel
methods, we will also evaluate it on the MIMIC dataset.

These evaluation and comparison strategies will help identify the most
effective model and setup for predicting treatment outcomes, enabling
clinical decision-makers to adopt the most reliable and interpretable
model for diabetes care.

\subsubsection{3.2.3 Feature Importance Analysis: Explainability and
Protocol
Refinement}\label{feature-importance-analysis-explainability-and-protocol-refinement}

Explainability is a critical requirement for medical informatics models,
as healthcare professionals must trust and understand predictions to
make informed decisions. To achieve this, feature importance analysis
will be conducted using \emph{SHAP} to identify the most influential
\emph{Quality Assessment Temporal Patterns} (\emph{QATPs}) associated
with patient outcomes. These patterns will help uncover key predictors
of complications or recovery, enhancing the interpretability of the
models.

Insights from \emph{SHAP} analysis will also guide protocol refinement
by highlighting the \emph{QATP}s most strongly correlated with clinical
outcomes. This ensures that the most impactful aspects of care are
prioritized in the updated protocol, ultimately improving clinical
decision-making and patient care (\textbf{Task No.4 in Section 3.2}).
The results will be validated by clinical experts to ensure their
alignment with established medical knowledge and their relevance to
real-world practice.

\subsubsection{3.2.4 Adherence Assessment Before and After
Interventions}\label{adherence-assessment-before-and-after-interventions}

In addition to feature importance analysis, this study will assess the
impact of two major interventions (\textbf{Task No.5 in Section 3.2}),
introduced by Dr. Irit Hochberg to the nursing teams, on guideline
adherence (\emph{GLs}) across six hospitals:

\begin{itemize}
\item
  The introduction of a diabetes GL learning software module for medical
  teams.
\item
  The implementation of a diabetes quality measure for managing
  hyperglycemia events.
\end{itemize}

To evaluate the effectiveness of these interventions, statistical
comparisons of \emph{GL}-based \emph{QA} scores will be conducted
between two distinct datasets:

\begin{itemize}
\item
  Pre-intervention phase: June 2022 to August 2023.
\item
  Post-intervention phase: August 2023 to January 2025.
\end{itemize}

This analysis will identify \emph{GLs} with significant implementation
challenges and evaluate how the interventions influenced adherence
patterns over time. Poorly adopted \emph{GLs} will be flagged for
further investigation, while their potential impact on patient outcomes
will also be assessed.

By combining these insights, the study aims to pinpoint barriers to
guideline adoption and provide actionable recommendations for improving
protocol compliance. These findings, when integrated with the feature
importance analysis, will support the development of a diabetes care
protocol that is not only effective but also practical and widely
adoptable in hospital settings.

\section{4 Initial Results and
Schedule}\label{initial-results-and-schedule}

\subsection{4.1 Initial Results}\label{initial-results}

This chapter documents what has been achieved to date and establishes
clear baselines for the remaining work. All metrics and curves below are
computed on the hold‑out validation cohort unless stated otherwise.

\begin{itemize}
\item
  \textbf{Knowledge Base (KB) Construction:\\
  } A comprehensive Knowledge Base for temporal abstractions has been
  created, encoding core clinical concepts and transformations (e.g.,
  states, gradients) in a structured, domain-consistent manner. The
  initial quality of these abstractions has been empirically tested
  using visualization tools and aggregation methods, as described in
  \textbf{{[}Chapter 3.2.1{]}} and shown in \textbf{{[}Appendix 1{]}},
  yielding promising preliminary plots and descriptive analyses.
\item
  \textbf{Embedder's Training:}
\end{itemize}

\begin{quote}
The embedder jointly optimises three auxiliary objectives
\,\textbf{BCE,\,\,MLM} and \,\textbf{Δt‑MSE}. Training completes in
\textasciitilde60 epochs with a smooth, monotonic loss decline down to a
total weighted loss \textasciitilde{} 0.1. This is a positive indication
for the ability of this architecture to represent the variance in the
data.

\includegraphics[width=4.31771in,height=2.91254in]{media/image5.png}

BCE Loss = 0.0199,

MLM Cross Entropy Loss = 0.806,

Δt MSE \textless{} 1e-4
\end{quote}

\begin{itemize}
\item
  \textbf{Transformer's Training:\\
  } An initial Transformer-based model has been trained on the
  structured event sequences, incorporating patient-specific context,
  multi-level event embeddings, and temporal features. The model showed
  a promising loss decrease during training:
\end{itemize}

\begin{quote}
\includegraphics[width=5.59573in,height=3.6095in]{media/image21.png}

This architecture remains experimental and requires further refinement
in how penalties are gradually increased during training. At present,
penalties rise during the warm-up phase (first five epochs) to prevent
them from dominating the loss too early. Under the current
configuration, this has led to underrepresentation of complication
logits, indicating that additional iterations are needed to ensure these
outputs receive sufficient emphasis during training. Early stopping was
triggered prematurely, suggesting that the current schedule does not yet
balance learning across all targets.

Focal BCE Loss = 0.0006,

Structural Penalties =0.0241,

Δt MSE = 0.0014
\end{quote}

\begin{itemize}
\item
  \textbf{Transformer's Classification Evaluation:\\
  } As described, the model was evaluated on 3 increasingly difficult
  tasks. While the first set of evaluation results has shown relatively
  low predictive performance, this phase has confirmed the viability of
  the architecture in handling complex, time-dependent event data and
  the practical challenges of working with real-world clinical data.
\item
  \textbf{Evaluation Framework and Levels of Analysis:\\
  } The initial Transformer evaluation was performed at three levels:

  \begin{itemize}
  \item
    \textbf{3-class multilabel outcome prediction} (e.g., release,
    complication, death)
  \end{itemize}
\end{itemize}

\begin{quote}
\includegraphics[width=2.89583in,height=1.17708in]{media/image10.png}

\includegraphics[width=5.55729in,height=0.44792in]{media/image20.png}

These results are still low and poorly represent the balance between the
3 classes.
\end{quote}

\begin{itemize}
\item
  \textbf{Multilabel prediction of specific clinical events}
\end{itemize}

\begin{quote}
Will be re-calculated after the model is fixed.
\end{quote}

\begin{itemize}
\item
  \textbf{Time-aware evaluation} of outcome timing relative to the
  admission trajectory (with 48 hours forgiveness window).
\end{itemize}

\begin{quote}
Will be re-calculated after the model is fixed.

These frameworks establish a baseline and highlight key limitations to
guide subsequent refinements.
\end{quote}

After finalizing the architecture, the next steps involve expanding the
knowledge base to include \textbf{patterns} and \textbf{pattern-level
adherence scores (based on QATPs)} and to assess the contribution of
guideline compliance signals to prediction accuracy.

\subsection{4.2 Planned Evaluation
Methods}\label{planned-evaluation-methods}

For the upcoming stages, we plan to evaluate the following aspects:

\begin{itemize}
\item
  \textbf{Expansion of the Knowledge Base (KB) and Pattern-Level
  Adherence:\\
  } The next steps involve expanding the knowledge base to include
  \textbf{temporal patterns} and their associated \textbf{pattern-level
  adherence (QA fuzzy scores)}. We will assess the contribution of these
  guideline compliance signals to overall prediction accuracy.
\item
  \textbf{Evaluation of Temporal Abstractions:\\
  } Using plots and descriptive statistics, we will evaluate the
  \textbf{quality, consistency, and clinical validity} of the
  abstractions - particularly for the patterns - across multiple patient
  scenarios.
\item
  \textbf{Comparative Evaluation:\\
  } We will compare the \textbf{Transformer-based architecture} with the
  \textbf{baseline TIRP/ TPF approaches} using predefined temporal
  features on evaluation levels 1 and 2 as described above. Furthermore,
  we will explicitly evaluate the impact of including \textbf{QA scores
  and guideline compliance signals} in the Transformer's training -
  comparing models with and without these features.
\item
  \textbf{Explainability Analysis:\\
  } We will conduct a SHAP-based feature importance analysis to identify
  which temporal patterns and QA scores contribute most to the
  prediction of clinical outcomes. This will help interpret the model's
  behavior and highlight critical protocol steps.
\item
  \textbf{Adherence Evaluation:\\
  } We will use standard statistical tests (T-Test) to evaluate the
  adherence level before and after the interventions, as described in
  \textbf{{[}Chapter 3.2.4{]}}. We will also try to identify guide-lines
  that were not well assimilated in the hospitals.
\end{itemize}

This comprehensive evaluation will ensure both the technical rigor and
the clinical relevance of our predictive models.

\subsection{4.3 Proposed Schedule}\label{proposed-schedule}

The following milestones outline the planned timeline for the next
phases of the research:

{\def\LTcaptype{none} % do not increment counter
\begin{longtable}[]{@{}
  >{\raggedright\arraybackslash}p{(\linewidth - 2\tabcolsep) * \real{0.5000}}
  >{\raggedright\arraybackslash}p{(\linewidth - 2\tabcolsep) * \real{0.5000}}@{}}
\toprule\noalign{}
\endhead
\bottomrule\noalign{}
\endlastfoot
\textbf{Milestone} & \textbf{Planned Completion} \\
\textbf{Data and Model Optimization}:

Finalize data cleaning and optimize the current dataset, addressing
inconsistencies and non informative records, and applying for more data
in case of need. Refine the Transformer model's architecture and
training setup based on initial performance insights. Finally, train and
evaluate the Transformer model \textbf{without} QA scores and
pattern-level features, establishing a baseline for performance. & July-
August 2025 \\
\textbf{Pattern and QATP Knowledge Base Development}:

Develop and integrate the knowledge base for temporal patterns and
pattern-level adherence into the data pipeline. Validate the generated
abstractions for clinical plausibility. & August- September 2025 \\
\textbf{Experiments: Transformer with QATPs and Adherence Scores}:

Train and evaluate the Transformer model \textbf{with} pattern-level
features and QA scores to assess their impact on predictive performance
and adherence insights. & September-

October 2025 \\
\textbf{Development and Evaluation of TPFs/TIRPs+VSM Models}:

Develop VSM-based predictive models using temporal pattern distributions
(TPF) and compare them to the Transformer models on eval levels 1 and 2
(3-class and multilabel outcome predictions). & November 2025 \\
\textbf{Adherence Evaluation Pre- and Post-Interventions}:

Conduct adherence analysis (using statistical tests such as T-Test) to
assess the effect of guideline-focused interventions and highlight
under-assimilated guidelines. & December 2025 \\
\textbf{Explainability and Feature Importance Analysis}:

Perform SHAP-based explainability analysis to highlight which features
and QA scores contribute most to predictions and identify protocol steps
with the strongest influence. & January 2026 \\
\textbf{Integration of Results and Chapter Drafting}:

Summarize findings across all experiments and write the results chapters
(initial results, evaluation, adherence impact). & February -- March
2026 \\
\textbf{Final Thesis Writing and Revisions}:

Prepare the final thesis document, incorporating revisions and final
edits based on feedback. & April 2026 \\
\textbf{Thesis Submission and Defense}:

Submit the final thesis and prepare for the defense. & May -- June
2026 \\
\end{longtable}
}

\section{References}\label{references}

\begin{enumerate}
\def\labelenumi{\arabic{enumi}.}
\item
  \hl{B. Ben Shahar, Y. Shahar, S. Jaffe, O. Cohen, E. Shalom, M.
  Selivanova, E. Rimon, A. Goldstein. "Supporting the Episodic
  Application of Clinical Guidelines over Significant Time Periods."
  \emph{Digital Health and Informatics Innovations for Sustainable
  Health Care Systems}. IOS Press, 2024. 1873-1877.‏}
\item
  \hl{E. Shalom, A. Goldstein, R. Weiss, M. Selivanova, N.M. Cohen, Y.
  Shahar. "Implementation and evaluation of a system for assessment of
  the quality of long-term management of patients at a geriatric
  hospital." \emph{Journal of Biomedical Informatics} 156 (2024):
  104686.‏}
\item
  \hl{Miksch, Silvia, Yuval Shahar, and Peter Johnson. "Asbru: a
  task-specific, intention-based, and time-oriented language for
  representing skeletal plans." \emph{Proceedings of the 7th workshop on
  knowledge engineering: methods \& languages (KEML-97)}. Milton Keynes,
  UK, The Open University, Milton Keynes, UK, 1997.‏}
\item
  \hl{Shahar, Yuval, and Mark A. Musen. "Knowledge-based temporal
  abstraction in clinical domains." \emph{Artificial intelligence in
  medicine} 8.3 (1996): 267-298.‏}
\item
  \hl{E. Shalom, Y. Shahar, Y. Parmet, E. Lunenfeld. "A
  multiple-scenario assessment of the effect of a continuous-care,
  guideline-based decision support system on clinicians' compliance to
  clinical guidelines." \emph{International journal of medical
  informatics} 84.4 (2015): 248-262.‏}
\item
  \hl{Shalom, Erez, Yuval Shahar, and Eitan Lunenfeld. "An architecture
  for a continuous, user-driven, and data-driven application of clinical
  guidelines and its evaluation." \emph{Journal of biomedical
  informatics} 59 (2016): 130-148.‏}
\item
  \hl{Peleg, Mor, Yuval Shahar, and Silvana Quaglini. "MobiGuide:
  guiding clinicians and chronic patients anytime, anywhere."
  \emph{Communications of the ACM} 65.4 (2022): 74-79.‏}
\item
  \hl{Chakravarty, Shubha, and Yuval Shahar. "CAPSUL: A constraint-based
  specification of repeating patterns in time-oriented data." Annals of
  mathematics and artificial intelligence 30 (2000): 3-22.‏}
\item
  \hl{Moskovitch, Robert, and Yuval Shahar. "Classification of
  multivariate time series via temporal abstraction and time intervals
  mining." Knowledge and Information Systems 45 (2015): 35-74.‏}
\item
  \hl{E.Sheetrit, N. Nissim, D. Klimov, Y. Shahar‏. "Temporal
  probabilistic profiles for sepsis prediction in the ICU." Proceedings
  of the 25th ACM SIGKDD International Conference on Knowledge Discovery
  \& Data Mining. 2019.‏}
\item
  \hl{Shknevsky, Alexander, Yuval Shahar, and Robert Moskovitch. "The
  semantic adjacency criterion in time intervals mining." Big Data and
  Cognitive Computing 7.4 (2023): 173.‏}
\item
  \hl{Raymond, Michal Weisman, and Yuval Shahar. "Provision of Decision
  Support Through Continuous Prediction of Recurring Clinical Actions."
  ICIMTH. 2023.‏}
\item
  \hl{Ashish Vaswani, Noam Shazeer, Niki Parmar, Jakob Uszkoreit, Llion
  Jones, Aidan N Gomez, Łukasz Kaiser, Illia Polosukhin. "Attention is
  all you need." Advances in neural information processing systems 30
  (2017).‏‏}
\item
  \hl{Pawel Renc, Yugang Jia, Anthony E. Samir, Jaroslaw Was, Quanzheng
  Li, David W. Bates \& Arkadiusz Sitek. "Zero shot health trajectory
  prediction using transformer." NPJ Digital Medicine 7.1 (2024): 256.‏}
\item
  \hl{Z. Yang, A. Mitra, W. Liu, D. Berlowitz, H. Yu‏. "TransformEHR:
  transformer-based encoder-decoder generative model to enhance
  prediction of disease outcomes using electronic health records."
  Nature communications 14.1 (2023): 7857.‏}
\item
  \hl{Lundberg, Scott M., and Su-In Lee. "A unified approach to
  interpreting model predictions." Advances in neural information
  processing systems 30 (2017).‏}
\item
  \hl{van de Beld, J. J., Pathak, S., Geerdink, J., Hegeman, J. H., \&
  Seifert, C. (2024, July). Feature importance to explain multimodal
  prediction models. A clinical use case. In World Conference on
  Explainable Artificial Intelligence (pp. 84-101). Cham: Springer
  Nature Switzerland.‏}
\item
  \hl{Lion, Matan, and Yuval Shahar. "Implementation and evaluation of a
  multivariate abstraction-based, interval-based dynamic time-warping
  method as a similarity measure for longitudinal medical records."
  Journal of Biomedical Informatics 123 (2021): 103919.‏}
\item
  \hl{Seyed Mehran Kazemi, Rishab Goel, Sepehr Eghbali, Janahan Ramanan,
  Jaspreet Sahota, Sanjay Thakur, Stella Wu, Cathal Smyth, Pascal
  Poupart, Marcus Brubaker. "Time2vec: Learning a vector representation
  of time." arXiv preprint arXiv:1907.05321 (2019)‏}
\item
  \hl{Shapiro, Michael, and Yuval Shahar. "Treatment Prediction in the
  ICU Using a Partitioned, Sequential, Deep Time Series Analysis."
  MEDINFO 2023 - The Future Is Accessible. IOS Press, 2024. 710-714.}
\item
  \hl{Akiba, T., Sano, S., Yanase, T., Ohta, T., \& Koyama, M. (2019,
  July). Optuna: A next-generation hyperparameter optimization
  framework. In Proceedings of the 25th ACM SIGKDD international
  conference on knowledge discovery \& data mining (pp. 2623-2631).‏}
\item
  \hl{Ribeiro, Marco Tulio, Sameer Singh, and Carlos Guestrin. "" Why
  should I trust you?" Explaining the predictions of any classifier."
  Proceedings of the 22nd ACM SIGKDD international conference on
  knowledge discovery and data mining. 2016.}
\item
  \hl{Shahar, Yuval, Silvia Miksch, and Peter Johnson. "The Asgaard
  project: a task-specific framework for the application and critiquing
  of time-oriented clinical guidelines." Artificial intelligence in
  medicine 14.1-2 (1998): 29-51.‏}
\item
  \hl{Shahar, Yuval. "A framework for knowledge-based temporal
  abstraction." Artificial intelligence 90.1-2 (1997): 79-133.‏}
\item
  \hl{IDF Atlas 10th Edition 2021.; 2021. https://diabetesatlas.org/}
\item
  \hl{Israel National Center for Disease Control Diabetes Registry
  Report 2012-2022.}
\end{enumerate}

\begin{quote}
\href{https://www.gov.il/BlobFolder/generalpage/diabetes-2012-2022/he/files_publications_units_ICDC_diabetes-2012-2022.pdf}{\hl{\ul{https://www.gov.il/BlobFolder/generalpage/diabetes-2012-2022/he/files\_publications\_units\_ICDC\_diabetes-2012-2022.pdf}}}
\end{quote}

\begin{enumerate}
\def\labelenumi{\arabic{enumi}.}
\setcounter{enumi}{26}
\item
  \hl{Summary of Cost of Diabetes in Israel 2018. National Health System
  and Health Policy Institution; 2018.
  \mbox{\href{https://www.btl.gov.il/Mediniyut/BakashatNetunim/dohot/Documents/Cost_of_Diabetes.pdf}{\ul{https://www.btl.gov.il/Mediniyut/BakashatNetunim/dohot/Documents/Cost\_of\_Diabetes.pdf}}}}
\item
  \hl{A. Radford, J. Wu, R. Child, D. Luan, D. Amodei, I. Sutskever.
  "Language models are unsupervised multitask learners." OpenAI blog 1.8
  (2019): 9.‏}
\item
  \hl{Karpathy, A. (2023). nanoGPT: The simplest, fastest repository for
  training/finetuning medium-sized GPTs. GitHub.
  \mbox{\href{https://github.com/karpathy/nanoGPT}{\ul{https://github.com/karpathy/nanoGPT}}}}
\item
  \hl{Devlin, Jacob, et al. "Bert: Pre-training of deep bidirectional
  transformers for language understanding." Proceedings of the 2019
  conference of the North American chapter of the association for
  computational linguistics: human language technologies, volume 1 (long
  and short papers). 2019.‏}
\item
  \hl{Paszke, Adam, et al. "Pytorch: An imperative style,
  high-performance deep learning library." Advances in neural
  information processing systems 32 (2019).‏}
\item
  \hl{Oded, S. (2025). ``Transform-EMR''. GitHub.
  \mbox{\href{https://github.com/shaharoded/Transform-EMR}{\ul{https://github.com/shaharoded/Transform-EMR}}}}
\item
  \hl{Oded, S. (2025). ``TAKAutomator''. GitHub.
  \mbox{\href{https://github.com/shaharoded/TAKAutomator}{\ul{https://github.com/shaharoded/TAKAutomator}}}}
\item
  \hl{Rappoport, Nadav, et al. "Kineret: Israel's Largest Hospital
  Network Transformed into the OMOP common data model for collaborative
  research." PloS one 20.10 (2025): e0334848.‏}
\end{enumerate}

\section{Appendices}\label{appendices}

\subsection{Appendix 1: Abstractions Quality
Assessment}\label{appendix-1-abstractions-quality-assessment}

\hl{As described in \textbf{{[}Chapter 4.1{]}}, a comprehensive
knowledge base for temporal abstractions was developed, encoding core
clinical concepts and transformations (e.g., states, gradients) in a
structured, domain-consistent manner. The initial quality of these
abstractions was assessed empirically using visualization tools and
aggregation methods to verify that the Mediator applies the intended
discrete thresholds and constructs intervals from raw records.}

\textbf{States thresholds verification:}

Plots were generated showing the occurrence distribution of every
discrete value for each abstracted measure. The goal was to confirm that
all defined thresholds are being used and that no concept shows an
implausible concentration of values (e.g., all cases in a single
category) or complete absence of certain categories. The results confirm
that thresholds are well-aligned with the data and produce a balanced
distribution across discrete states.

\includegraphics[width=6.27083in,height=2.89276in]{media/image28.png}

\includegraphics[width=6.27083in,height=2.88021in]{media/image16.png}

\includegraphics[width=6.26772in,height=3.16667in]{media/image8.png}\includegraphics[width=6.26772in,height=3.20833in]{media/image26.png}

\includegraphics[width=6.27083in,height=3.03646in]{media/image2.png}

\includegraphics[width=6.26772in,height=3.33333in]{media/image9.png}

\includegraphics[width=6.26772in,height=3.25in]{media/image27.png}

\includegraphics[width=6.27083in,height=2.81771in]{media/image6.png}

\textbf{\hl{Trends parameters verification:}}

\hl{Trend abstractions require the definition of both a minimum
significant change and a time window in which that change must occur. If
increase or decrease states are absent from the data, it suggests the
parameters were misconfigured or unused. This QA step checked for the
presence of both increase and decrease trends across all measures. The
findings confirm that trend parameters were applied correctly, and the
abstractions capture meaningful temporal changes.}

\includegraphics[width=6.26772in,height=3.375in]{media/image14.png}

\includegraphics[width=6.26772in,height=3.33333in]{media/image23.png}

\includegraphics[width=6.26772in,height=3.31944in]{media/image3.png}

\includegraphics[width=6.27083in,height=3.11775in]{media/image25.png}

\includegraphics[width=6.27083in,height=3.03646in]{media/image22.png}
\hl{}

\textbf{\hl{Intervals creation verification:}}

\hl{This check evaluated whether the relevance duration for each measure
or medication is appropriate given its appearance frequency in the EMR.
The Mediator's objective is to merge raw timestamped records into
clinically meaningful intervals. If the process results only in discrete
points without interval formation, temporal continuity is lost. For each
concept, the average duration of abstracted intervals and the ratio of
abstracted to raw records were calculated. In most cases, horizontal
support decreased substantially after abstraction, indicating successful
interval creation. Concepts with little or no change generally appeared
only once per patient, making longer interval abstraction impossible.}

\includegraphics[width=6.26772in,height=3.44444in]{media/image18.png}

\includegraphics[width=6.26772in,height=3.63889in]{media/image15.png}

\includegraphics[width=6.26772in,height=3.36111in]{media/image13.png}

\includegraphics[width=6.26772in,height=0.97222in]{media/image4.png}

\subsection{}\label{section-1}

\subsection{Appendix 2: Automated TAK Generation
Tool}\label{appendix-2-automated-tak-generation-tool}

\hl{As described in \textbf{{[}Chapter 3.1.1{]}}, the Mediator depends
on a comprehensive library of TAK (Temporal Abstraction Knowledge) XML
files, each encoding a discrete clinical abstraction such as a
\emph{state}, \emph{trend}, \emph{context}, or \emph{event}. While these
files follow a strict XML schema, they must also conform to business
logic parameters defined in the project's Knowledge Base (KB)
spreadsheet. Maintaining and updating hundreds of such files manually is
error-prone, especially when thresholds, labels, or logic rules need to
be adjusted globally.}

\hl{To address this challenge, the \textbf{TAKAutomator} {[}33{]} tool
was developed. This Python-based framework automates the creation and
curation of TAK files from a structured Excel specification, ensuring
both structural and semantic consistency. The tool allows the Knowledge
Base to be treated as a \emph{single source of truth}, from which all
TAK files can be generated and re-generated at any time.}

\textbf{Method}

\hl{The TAKAutomator operates in three main stages:}

\begin{enumerate}
\def\labelenumi{\arabic{enumi}.}
\item
  \hl{\textbf{Excel Parsing\\
  } The tool ingests a preformatted Excel workbook containing
  definitions for all TAK entities. Each sheet represents a TAK type
  (e.g., \emph{states}, \emph{trends}, \emph{contexts}) and includes
  IDs, parameter thresholds, labels, mappings, and any other required
  metadata. The first step when running the program is an automated
  validation of the business logic represented in this excel:}

  \begin{itemize}
  \item
    \hl{Dependencies on existing raw concepts.}
  \item
    \hl{Proper range coverage in states.}
  \item
    \hl{Non conflicting values between entities.}
  \end{itemize}
\item
  \hl{\textbf{Template Rendering\\
  } For each row in the Excel file, TAKAutomator retrieves the
  appropriate XML template, substitutes placeholders with the
  corresponding Excel values, uses LLM for more complex generations
  (treating this as a \emph{Retrieval-Augmented Generation (RAG)} task)
  and generates a TAK file. This step ensures consistent formatting and
  eliminates manual XML editing.}
\item
  \hl{\textbf{Validation Loop\\
  } Every generated TAK file is passed through a Validation module,
  which checks:}

  \begin{itemize}
  \item
    \hl{Structural validity against the XML schema.}
  \item
    \hl{Compliance with business logic rules from the Excel KB.}
  \item
    \hl{Presence and correctness of thresholds, allowed values, and
    mappings.\\
    The process can be iterated to automatically correct fixable
    issues.\\
    }
  \end{itemize}
\end{enumerate}

\textbf{Contribution}

\hl{The TAKAutomator has significantly accelerated the curation and
maintenance of the \emph{Mediator}'s Knowledge Base. Using it, a
repository of \textbf{200+ TAK files} was created and is continuously
updated in sync with the Excel definitions. Global changes---such as
updating state thresholds or adjusting trend definitions-can now be
propagated to the entire library with a single command, eliminating
repetitive manual edits.}

\hl{This automation ensures:}

\begin{itemize}
\item
  \hl{\textbf{Consistency} between the Knowledge Base and the TAK
  files.}
\item
  \hl{\textbf{Scalability}, as new abstractions can be added with
  minimal overhead.}
\item
  \hl{\textbf{Auditability}, since changes to the Excel KB automatically
  produce a new validated TAK set.}
\item
  \hl{\textbf{Explainability}, by enabling medical experts-without
  programming expertise-to independently update and refine the KB.
  Through the structured Excel interface, they can define states,
  trends, and thresholds in clinically meaningful terms, which are then
  automatically translated into validated TAK files for use by the
  \emph{Mediator}.}
\end{itemize}

\hl{By integrating TAKAutomator into the Mediator pipeline, the lab has
gained a sustainable and error-resistant process for managing temporal
abstraction definitions-transforming what was once a manual,
high-friction task into a streamlined and reproducible workflow.}

\end{document}
